\documentclass[11pt,letterpaper]{article}
\usepackage[T1]{fontenc}
\usepackage[utf8]{inputenc}
\usepackage{lmodern,amsmath,amssymb,amsthm}
\usepackage[margin=1in]{geometry}
\usepackage{microtype,booktabs,longtable,array,xcolor,listings,enumitem}
\usepackage[hidelinks]{hyperref}
\usepackage{fancyhdr}
\setlength{\headheight}{14pt}
\pagestyle{fancy}
\fancyhf{}
\fancyhead[L]{\small Deterministic specialization for iPRA}
\fancyhead[R]{\small Research report}
\fancyfoot[C]{\small\thepage}
\setlength{\parindent}{0pt}
\setlength{\parskip}{4pt}
\setlist{itemsep=2pt,topsep=5pt}
\emergencystretch=2em
\newtheorem{theorem}{Theorem}[section]
\newtheorem{lemma}[theorem]{Lemma}
\newtheorem{proposition}[theorem]{Proposition}
\newtheorem{corollary}[theorem]{Corollary}
\theoremstyle{definition}
\newtheorem{definition}[theorem]{Definition}
\newtheorem{remark}[theorem]{Remark}
\lstset{basicstyle=\ttfamily\scriptsize,columns=fullflexible,
 keepspaces=true,breaklines=true,showstringspaces=false,
 frame=single,rulecolor=\color{black!20},xleftmargin=4pt,xrightmargin=4pt}
\hypersetup{pdftitle={Deterministic specialization for iPRA},
 pdfauthor={Research prepared for Vladimir Reshetnikov},
 pdfsubject={A complete five-rule presentation with no numerical target in clause application}}
\begin{document}
\begin{titlepage}
\vspace*{0.7cm}
{\large Mathematical foundations\hfill 11 November 2026\par}
\vspace{1.0cm}
{\LARGE\bfseries Deterministic specialization\\[5pt]for iPRA\par}
\vspace{0.6cm}
{\Large Removing the term argument\\from clause application\par}
\vspace{0.8cm}
{\large Prepared for Vladimir Reshetnikov\par}
\vspace{0.8cm}
\begin{abstract}
The current ConciseIPRA proposals already reduce full intuitionistic
primitive recursive arithmetic to five rule patterns on equations and
one paired binder. We make a further inferential restriction: a clause
application supplies no numerical target. The newest clause is always
specialized at one, and every older clause at the newest eigenparameter.
All arbitrary numerical specialization remains derivable.

The completeness proof names the intended term by an existential
equality. Its main difficulty is structural: extending a scope changes
the permitted target of an existing inference. A rank-recursive transport
construction and a proof-height compiler regenerate premises under the
larger scope; unrestricted weakening is not assumed. We prove exact
consequence equivalence with the existing clause calculus, preserving its
whole-domain bi-interpretation and equiconsistency with full first-order
iPRA. Both numerical interpretation composites have graph \(x=y\).

The complete readable checker has 2914 normalized bytes, 819 significant
tokens, and 896 AST nodes, versus 2950, 829, and 902 for the incoming
uniform-Step checker. We separate the small rule-interface reduction
from an independent elimination of redundant scope state. A separate
implementation of the unchanged native Eq/Q/D calculus reduces its
fully validated checker from 866 to 820 bytes. These are controlled
improvements, not a proof of an absolute minimum. The package includes
complete written arguments, a bidirectional executable translation,
100 forward and reverse certificate checks, and independent boundary
audits. It does not claim a new Lean or Rocq formalization.
\end{abstract}
\vfill
\textbf{Frozen source.} ProveIt, commit\\
\texttt{5d85023b0688e03b75f0ec3a04bb43ef0bb51705}.\\
\textbf{Main change.} An arbitrary-term inference becomes deterministic.\\
\textbf{Tradeoff.} Smaller checking interface; potentially much larger proofs.
\end{titlepage}
{\small\setlength{\parskip}{0pt}\tableofcontents}
\clearpage

\section{Result, baseline, and precise claim}

The question asks for a presentation simpler than the current
\texttt{ConciseIPRA} proposal, while
preserving full iPRA up to bi-interpretation. We use the complete checker
as the implementation unit and include all syntax validation, scopes,
normalization, inference checking, JSON input, and command-line output.
The comparison is pinned to the commit on the title page, not to a moving
branch~\cite{baseline}.

The repository index describes fifty-six filed research packages. Its
five incoming archives were also inspected: computation normal form,
counter evaluation, reduced clause calculus, root-free clausal logic,
and uniform successor Step. In particular, the five-rule architecture,
zero-normal iterators, Refl plus uniform Step, and the clause interface
were already available~\cite{uniform,rootfree,reduced}. Combining those
ingredients alone would not be a new mathematical reduction.

The present change concerns the remaining arbitrary specialization
argument. An unrestricted clause instruction has the shape
\[
  [\texttt{u},\ i,\ t,\ \pi],
\]
where i chooses a retained clause, t is any numerical term, and
\(\pi\) proves its antecedent. Our instruction is
\[
  \boxed{[\texttt{u},\ i,\ \pi]}.
\]
The permitted target is computed from the clause position and the current
parameter scope. No hidden arbitrary term, mode bit, substitution
certificate, or evaluator is moved to another field of this instruction.
The existing existential introduction still supplies a witness; the
completeness translation uses that existing facility to name the desired
term.

\begin{theorem}[Main result, stated in advance]
The deterministic five-rule calculus of
Section~\ref{sec:deterministic-specialization} has exactly the same
consequence relation as the unrestricted zero-normal clause calculus.
Its ordinary realization is bi-interpretable and equiconsistent with
full first-order intuitionistic primitive recursive arithmetic. Both
interpretations use the whole numerical domain and actual equality;
the internally definable isomorphisms for the numerical composites
are \(x=y\).
\end{theorem}

The mathematical improvement is a restriction of a primitive rule, with
an effective recovery of every omitted instance. The rule count stays
five. It is not a deletion of induction, equality, or a logical
connective. There is also no theorem here that the resulting presentation
is globally simplest among all effective presentations. Such a statement
would require a fixed representation and a much stronger lower-bound
argument than a matched implementation comparison.

\subsection{What is counted and what is inherited}

The selected numerical language, paired connective, iterator arithmetic,
and scalar-to-iPRA bridge are inherited. We reproduce their essential
arguments and the five authenticated bridge sources in the appendices.
The new contribution is deterministic specialization, its persistence
proof, and the executable compiler. The separate grammar/scope traversal
fusion is an implementation refinement of the old native calculus.

These contributions should not be conflated. The 2914-byte readable
checker implements the new deterministic clause calculus. The 820-byte
supplement implements the unchanged, differently encoded native Eq/Q/D
calculus. It is not a 820-byte implementation of the new calculus.
All measured programs import only the Python standard library, and
their complete sources are supplied.

\section{The complete deterministic-specialization calculus}
\label{sec:deterministic-specialization}

\subsection{Syntax, computation, and all primitive rules}

There is one numerical sort.  Start from finite raw terms
\[
 t ::= x\mid 0\mid S(t)\mid I(n,a;u.h).
\]
Only u in h is bound by the iterator; n and a are outside that binder.
The zero-normalization map acts recursively on subterms and contracts
\(I(0,a;u.h)\) to a, but does not contract successor-counter iterators.
The terms of the present calculus are the fixed points of this map.
Every substitution is capture avoiding and followed by this normalization,
including a substitution into an iterator counter that newly becomes zero.
For readability, expressions below always denote their zero-normal forms.

Formulas and clauses have grammars
\[
 F ::= (s=t)\mid D_{x,y}F,\qquad [A;y.B(y)]\quad(y\notin\mathrm{FV}(A)).
\]
Both x and y are bound by D.  A context is a finite ordered list of clause
occurrences; occurrences remain distinct even when their contents agree.
An ordinary assumption A is represented by the constant clause
\([0=0;y.A]\), with y unused.  Put \(1=S0\), and designate \(0=1\)
as contradiction.  Clause order is used only to choose a specialization
target; clauses are reusable, and no assumption is consumed by a rule.

The following five schemas are the entire primitive inventory.  There
are no additional logical, equality, computation, successor, induction,
weakening, exchange, or cut rules.
\[
 \frac{}{\Gamma\vdash t=t}\ (\mathsf{Refl}),\qquad
 \frac{\Gamma\vdash I(c,a;u.h)=t\qquad
       \Gamma\vdash I(Sc,a;u.h)=r}
      {\Gamma\vdash r=h[t/u]}\ (\mathsf{Step}).
\]
In Step the first iterator is normalized, so it is just a when c is zero.
Repeated expressions must be syntactically identical; no hidden equality
decision or evaluation oracle occurs in this rule.
\[
 \frac{\Gamma\vdash D_{x,y}B(x,y)\qquad
       [B(b,0);y.B(b,y)],\Gamma\vdash F}
      {\Gamma\vdash F}\ (\mathsf{DE}).
\]
Here b is fresh for \(\Gamma,F\) and all passive matrix parameters.
\[
 \frac{[0=0;v.B(t,\eta)],\Gamma\vdash B(t,Sz)}
      {\Gamma\vdash D_{x,y}B(x,y)}\ (\mathsf{DI}),\qquad
 \eta=\begin{cases}
 z,&B\text{ is a literal equation and }t\text{ is literally }0,\\
 0,&\text{otherwise}.
 \end{cases}
\]
Here v is unused, and z is fresh for \(\Gamma,t\) and all passive matrix
parameters.  An equation beneath a D is not a literal equation.  The
fifth schema is Use$_*$ below.  The witness-selected DI convention and
its recovery of the missing equation-case instances at witness zero are
inherited from the existing uniform-Step calculus; they are not altered
by the present specialization reduction.

Store context
clauses newest first.  For a nonempty parameter scope let \(b_*\) be its
last parameter; for the empty scope put \(b_*=0\).  Replace unrestricted
clause use by the following rule:
\[
 \frac{\Gamma\vdash A}{\Gamma\vdash B(v_i)}
 \quad [A;y.B(y)]=\Gamma_i,
 \qquad v_i=\begin{cases}1&i=0,\\ b_*&i>0.\end{cases}
 \tag{Use$_*$}
\]
Thus the newest clause is used at \(1=S0\); every older clause is used
at the current last parameter.  The certificate supplies its clause
index and antecedent proof, with no numerical target or mode bit.

\subsection{The arithmetic transport used in the proof}

For \(J_s=I(1,s;u.u)\), Step sends premises \(s=t\) and \(J_s=r\)
to \(r=t\).  If \(p:s=t\) and \(q:s=r\), then
\[
 \mathsf{Step}\bigl(p,\mathsf{Step}(q,\mathsf{Refl}(J_s))\bigr)
       :r=t.
 \tag{Pivot}
\]
This gives symmetry by putting the first premise equal to \(s=s\), and
then transitivity.  For a term template h, put \(K=I(1,s;u.h)\).
Step on reflexivities gives \(K=h(s)\); Step on \(p:s=t\) and that
equation gives \(h(s)=h(t)\).  Hence transport in an equation template
\(F(v)=(a(v)=b(v))\) follows from
\[
 a(t)=a(s),\qquad a(s)=b(s),\qquad b(s)=b(t).
\]
This derivation uses only Refl and Step and opens no new scope.
Applying Step to the two appropriate reflexivities also gives successor
computation.  In particular \(C(0,a,b)=a\) is definitional and
\(C(Sz,a,b)=b\) is derivable for the selector
\(C(y,a,b)=I(y,a;u.b)\).  These observations prove every arithmetic
operation needed by the specialization construction independently of
the logical rules.

\subsection{Exact consequence equivalence}

\begin{theorem}[Deterministic-specialization equivalence]
The resulting calculus derives exactly the same open sequents as the
unrestricted five-rule clause calculus.  The ordinary numerical and
formula interpretations are unchanged.
\end{theorem}

\begin{proof}
Soundness is immediate because every new use is an old use.  We give a
constructive completeness proof.  Scope extension is not assumed: it
changes the last parameter, and inserting a clause changes which clause
is newest.  Those two effects are handled by the proof-producing
procedures below.

Let \(C(y,a,b)=I(y,a;u.b)\), and let
\[
 N_t=D_{x,y}\bigl(C(y,0,x)=C(y,0,t)\bigr).
 \tag{Name}
\]
Introduce \(N_t\) at witness \(t\).  Its required successor equation
has literally equal sides, so reflexivity proves the branch in either
introduction mode.  Eliminate \(N_t\) and call the fresh parameter
\(b\).  The newly inserted clause is
\[
 \mathcal N_{b,t}=[0=0;\ y.C(y,0,b)=C(y,0,t)].
 \tag{N}
\]
Since it is newest, Use$_*$ applies it at 1.  Reflexivity supplies its
antecedent, and selector computation gives \(b=t\).  This particular
proof is valid immediately after insertion; we do not pretend that it
automatically remains valid after further context extensions.

A \emph{persistent premise procedure} is a procedure that produces a
proof of its conclusion after every injective renaming of its free
parameters and every embedding of its retained clause occurrences into
a larger context and scope.  This definition describes an external
finite proof-producing algorithm, not an added structural inference
rule.  A retained constant clause for an ordinary assumption is
persistent: either permitted target gives the same formula.  Refl and
arithmetic combinations of persistent premises are also persistent.

\paragraph{Atomic specialization comes first.}
Let \([A;y.E(y)]\) be any retained clause whose body E is a literal
equation, let P be a persistent procedure for A, and let t be any term.
To specialize E at t, introduce and eliminate \(N_t\).  In its branch,
request P in the enlarged context, and apply the original clause at the
new last parameter b.  The original clause is now older, so this is
exactly its permitted use.  The result is \(E(b)\).  The immediately
available proof \(b=t\) from (N), together with equational congruence,
gives \(E(t)\).  Close the elimination.  Call this procedure \(U_0\).
The arithmetic transport here opens no scope, so it does not require a
persistent proof of \(b=t\).  The whole recipe is uniform under context
and scope embeddings, so \(U_0\) produces a persistent output whenever
P is persistent.

In particular, after (N) has ceased to be newest, its equation \(b=t\)
still has a persistent proof: use \(U_0\) to specialize that retained
atomic clause at 1, with Refl of zero as its persistent antecedent, and
apply selector computation.  This use of \(U_0\) is not circular.
Its own auxiliary naming clause is newest when used at 1, and its final
transport is atomic.  This is the essential extra step that permits the
deterministic protocol.

\paragraph{Mutual recursion by formula rank.}
Rank is the number of D binders above the final equation.  We construct
\(T_r\) and \(U_r\) for every finite rank r:
\begin{itemize}
\item \(T_r\) transports a formula F of rank r along a persistent proof
of \(s=t\), from an arbitrary supplied proof of \(F(s)\) to a proof of
\(F(t)\).  The supplied proof of \(F(s)\) need not be persistent.
\item \(U_r\) specializes a retained clause \([A;y.F(y)]\), with F of
rank r, at an arbitrary term t, provided its antecedent A is supplied by
a persistent procedure.  It produces a persistent proof procedure.
\end{itemize}
For \(T_0\), use the equational congruence already derived from Step;
\(U_0\) was constructed above.

Assume the operations are known below rank r, and write
\(F(v)=D_{x,y}G(v,x,y)\).  For \(T_r\), use the arbitrary supplied
proof of \(F(s)\) immediately as the major premise of DE.  Let its fresh
witness be c.  The context now contains
\([G(s,c,0);y.G(s,c,y)]\).  Introduce \(F(t)\) at witness c.
Because c is a fresh parameter, it is not literally zero; this is a case
introduction, including when G is an equation.  The local assumption
\(G(t,c,0)\) is a constant clause and therefore persistent.  Transport
it back along the persistent equality by \(T_{r-1}\), obtaining a
persistent proof of \(G(s,c,0)\).  Use \(U_{r-1}\) to specialize the
retained closure at \(Sz\), then transport forward by \(T_{r-1}\) to
obtain \(G(t,c,Sz)\).  Close DI and DE.  The original arbitrary minor
proof was used outside the new scope and was never weakened.  The recipe
is persistent when that minor premise is supplied persistently.

Now construct \(U_r\).  Introduce and eliminate \(N_t\), obtaining
the newest parameter b.  Request the antecedent procedure in that
extension and use the original, now older, clause at b.  This yields
\(F(b)\).  The equality \(b=t\) has the persistent proof from the
atomic-specialization paragraph, even under the scopes later opened by
\(T_r\).  Apply \(T_r\) to obtain \(F(t)\), and close the naming
elimination.  This construction is uniform under embeddings.  It uses
\(T_r\) and the already established \(U_0\); the construction of
\(T_r\) used only ranks below r.  Hence the recursion is well founded.

\paragraph{Compiling arbitrary derivations.}
By structural recursion on an unrestricted proof, produce a persistent
procedure for its concluding sequent.  Refl and Step use the procedures
for their immediate subproofs.  At an unrestricted Use, its antecedent
procedure is already available by recursion on the proper subproof;
invoke \(U_r\) at the supplied numerical target.  At DE or DI, choose
fresh parameters in the requested larger scope, extend the context
embedding by the new clause, and invoke the previously constructed
procedure for the proper branch subproof.  Injective parameter renaming
preserves the formula root and the literal-zero witness test, so it
preserves the introduction mode.  Scope freshness and the clause-index
mapping are explicit in this procedure.

This compilation recurses on original proof height.  Its auxiliary T/U
procedures recurse on formula rank as described above.  There is no
recursion on an already generated proof of uncontrolled size.  Every
finite input proof therefore produces a finite restricted proof.
Evaluate the persistent output at the identity context embedding to
obtain the required derivation.
\end{proof}

\paragraph{Implementation significance.}
The matched complete checker replaces
\texttt{[u, index, target, premise]} by \texttt{[u, index, premise]}, removes
target-term validation from Use, and reverses the context-storage order.
The selected target is computed from the index and current scope.
For the incoming tagged uniform-Step implementation, the complete
normalized source changes from 2950 to 2938 bytes.  Significant tokens
remain 829, AST nodes change from 902 to 903, and informational gzip-9
size changes from 1005 to 999 bytes.  This is a mathematical removal of
the arbitrary specialization argument, with a small byte reduction;
it is not a strict improvement in every size coordinate or an absolute
minimality result.

\subsection{Ordinary realization and the interpretation composites}

Use the unchanged ordinary readings
\[
 \bigl(D_{x,y}B\bigr)^\circ
   =\exists x\bigl(B^\circ(x,0)\to\forall y\,B^\circ(x,y)\bigr),
 \qquad
 [A;y.B]^\circ=A^\circ\to\forall y\,B^\circ.
\]
Equality is actual numerical equality.  The underlying ordinary theory is
intuitionistic first-order logic with equality, the scalar iterator
equations, numerical coverage
\(\forall x(x=0\lor\exists y\,x=Sy)\), separation \(0=S0\to\bot\),
and induction for literal equations with side assumptions.  The inherited
uniform-Step equivalence identifies this ordinary theory with the
unrestricted five-rule clause presentation.  Its logical macro theorem
recovers all intuitionistic connectives and both quantifiers; its scalar
arithmetic bridge derives predecessor, reflection of equality, complete
quantifier-free induction, pairing, and a compiler for every primitive
recursive definition.

The new equivalence changes no term, formula, clause reading, or
arithmetical axiom.  It supplies proof translations in both directions on
exactly the same syntax: one direction includes Use$_*$ as the indicated
ordinary Use, and the other is the finite construction in the theorem.
Consequently it composes with the inherited whole-domain bi-interpretation
with full first-order iPRA.  The numerical domains in both directions are
the whole domain, equality remains actual equality, and the numerical
composite isomorphisms have graph \(x=y\).  The formula composites are
the inherited internally provable equivalences of the logical macros;
the present change adds an identity translation on formulas.  Thus
bi-interpretability and equiconsistency follow without a new arithmetic
bootstrap.  This paragraph relies on the inherited bridge theorem; the
new proof above establishes the specialization equivalence itself.

\paragraph{Scope of the novelty claim.}
The new item is the restriction of clause specialization to a target
computed solely from clause position and current scope, together with
the persistent-procedure argument that restores arbitrary specialization
and structural scope extension.  The iterator, five-rule architecture,
clausal contexts, scalar bridge, and arithmetic equality bootstrap are
inherited.  The reviewed source snapshot and the summaries of its filed
packages, together with the incoming reduced-clause, root-free-clause,
and uniform-Step packages, did not contain this deterministic target rule.
This is a scoped comparison, not a claim of priority over all possible
axiomatizations or a proof of global minimality.


\section{Closing the bridge to the inherited axiomatization}
\label{sec:closing-bridge}

The main theorem is a proof transformation between two presentations
with identical formulas. This section records the remaining links to
the inherited Choice-DI calculus, so that the new restriction does not
rest on an unexamined claim that five rules happen to be sufficient.
All derivations in this section first take place in the unrestricted
clause calculus; the main compiler then transports them into the new one.

\subsection{Normalization and the arithmetic interface}

Zero normalization terminates: contracting \(I(0,a;u.h)\) deletes its
outer iterator and unused body. The recursively defined normalizer
first normalizes the children and then either returns the initializer
or rebuilds the iterator. This determines a unique normal form without
performing any nonzero iteration. Substitution is followed by that same
normalizer. Structural induction on a raw term proves that it equals
its normal form using the raw zero equation and congruence. Structural
induction on terms also proves coherence of normalizing substitution.

The Pivot derivation in the main theorem recovers all ordinary equality
reasoning from Refl and Step. Successor computation is Step on the
reflexivities of \(I(c,a;u.h)\) and \(I(Sc,a;u.h)\). Conversely,
the raw computation equation and equality transport prove Step:
replace \(I(c,a;u.h)\) by t in its update, and replace the successor
iterator by r. Thus no arithmetic strength is added by the uniform
evaluation rule. No general primitive-recursive evaluation algorithm
is part of normalization.

\subsection{All case introductions and all equation inductions}

The unified DI uses progress at an equation with witness literally zero.
The omitted case instance is nevertheless derivable. Let
\[
 c=I(1,0;u.u).
\]
It is a nonliteral zero-normal term, and Step proves \(c=0\). Given a
case premise \(\Gamma,E(0,0)\vdash E(0,Sz)\), begin instead with
\(E(c,0)\), transport to its zero-witness instance, insert the given
premise, and transport its result back to \(E(c,Sz)\). Introduction at
c is case mode and concludes the original \(D_{x,y}E(x,y)\). Formula
roots and the witness are explicit finite syntax. Ordinary-assumption
cut here is proved by induction on the unrestricted derivation: a use
of the distinguished constant clause is replaced by its supplied proof.

Progress at any fixed witness t is recovered by introducing
\(D_{w,y}E(t,y)\) with unused w and witness zero. Elimination supplies
\([E(t,0);y.E(t,y)]\). Under the desired original base, specialize it
at \(Sz\), then apply the derived case introduction at t. In
particular, from a base \(E(0)\) and a step under \(E(z)\), derive
\(E(t)\) for every numerical target t. Compound progress is never
introduced; that would change the induction strength.

Transport through a D formula follows by induction on its nesting
depth. Eliminate the old D, introduce the new one at the fresh witness,
transport its base backwards, specialize the old clause, and transport
its successor instance forwards. Every recursive transport acts on the
strictly smaller matrix. This gives the full formula-transport rule of
the inherited Choice-DI presentation. The reverse simulation uses its
arithmetic equations and transport to recover Refl and Step, and its
case or progress DI to recover our selected DI. The zero-normalized
and raw presentations therefore have the same translated sequents.

\subsection{Why this is full intuitionistic first-order arithmetic}

The primitive binder has ordinary reading
\[
 D_{x,y}B=\exists x\bigl(B(x,0)\to\forall y B(x,y)\bigr).
\]
Although the primitive formula grammar is small, it represents all the
intuitionistic connectives, including positive disjunction and existential
quantification. This is stronger than showing that an implicational or
negative fragment is mutually interpretable with arithmetic.

For reference, let \(G_vA\) be the equation-wise selector gate and
\(M_v(A,B)\) the formula selector constructed in
Appendix~\ref{subsec:choice-logic}. Their internally proved interfaces are
\[
 \vdash G_0A,\quad G_{Sv}A\dashv\vdash A,\quad
 M_0(A,B)\dashv\vdash A,\quad M_{Sv}(A,B)\dashv\vdash B.
\]
With \(p(z)=I(z,0;u.C(u,1,0))\), the macros are
\begin{align*}
 A\Rightarrow_*B&=D_{u,z}M_z(A,B),\\
 A\wedge_*B&=D_{u,z}G_zM_{p(z)}(A,B),\\
 \exists_*x\,A(x)&=D_{x,y}G_yA(x),\\
 A\vee_*B&=\exists_*x\,M_x(A,B),\\
 \forall_*x\,A(x)&=A(0)\wedge_*D_{u,z}G_zA(z),\\
 \bot_*&=(0=1).
\end{align*}
Unused and freshly introduced variables are chosen to avoid capture.
The appendix defines the gate and selector recursively even when the
operands have different D depths, and proves their inference rules and
both formula round trips. These macros belong to the external
translation, not to the trusted checker.

The ordinary realization uses actual equality, the two iterator
equations, numerical coverage, literal-equation induction with side
assumptions, and the explicit separation bridge
\((0=S0)\to\bot\). Stating this bridge matters: positive equations alone
admit a singleton structure. In the native language, contradiction is
the designated equation \(0=S0\), and its explosion rule is derived;
the separation bridge translates back to an identity implication.

The inherited internal arithmetic proof has a strict dependency order:
scalar predecessor; actual-equality reflection; characteristic equations
for all quantifier-free formulas and their induction; pairing and
projections; and the translation of arbitrary primitive recursion.
For a recursion with base g and step h, the final compiler stores
\((n,f(n))\) in a single scalar iterator state. Its projection equations
are proved internally for arbitrary elements, not checked only on
standard numerals. The full proofs appear in
Appendices~\ref{sec:hybrid-bootstrap}, \ref{sec:eic-completion}, and
\ref{sec:iterator-coding}.

In the opposite direction, each fixed iterator is a primitive-recursive
definition in iPRA. Internal equation induction proves uniqueness of
both iterator definitions and ordinary primitive recursions. Induction
on a finite function description consequently proves that its double
translation is the original operation. Every number remains in the
domain, equality remains equality, and the identity graph is the
internally provable numerical isomorphism. Formula composites use the
logical macro equivalences. This establishes bi-interpretability, not
merely agreement in the standard natural numbers or equiconsistency.

\section{The complete checker and certificate contract}
\label{sec:checker}

\subsection{Concrete syntax and scopes}

The primary checker is \path{code/checker.py}. It is one standalone
module, with no custom imports or callbacks to the proof generator.
The input is exactly a decoded JSON object with keys \texttt{goal}
and \texttt{proof}. The goal must be a closed native formula. Success
prints \texttt{ACCEPT} and exits with status zero; every caught invalid
certificate exits with a nonzero status. The implementation does not
use Python assertions for acceptance.

\begin{center}
\begin{tabular}{ll}
\toprule
Object & JSON representation\\
\midrule
Zero & \texttt{0}\\
Successor & \texttt{['s', t]}\\
Nonzero-counter iterator & \texttt{['i', n, a, h]}\\
Equation & \texttt{['=', s, t]}\\
Paired formula & \texttt{['d', B]}\\
Bound numerical variable & Positive integer index\\
Free eigenparameter & Negative integer level\\
\bottomrule
\end{tabular}
\end{center}
Single quotes in this notation are explanatory Python notation; actual
JSON files use double quotes. Zero-counter iterator nodes are rejected
in supplied syntax. The opening function contracts them if substitution
creates them. An iterator adds one binder only in its body. A D matrix
adds two: index 1 is its y coordinate and index 2 its x coordinate.
Positive indices beneath additional binders are adjusted accordingly.

At scope k the permitted free levels are \(-k,\ldots,-1\).
The newest is \(-k\), and the next is \(-k-1\). Every supplied term
operand is locally closed with respect to positive indices; its free
variables are the negative levels. This invariant is why opening can
insert the supplied value beneath a binder without an additional shift
of that value. Opening still tracks the binder depth, removes the selected
index, and decrements more distant bound indices. Tests include values
occurring under nested iterator and D binders.

The proof shapes are:
\begin{center}
\begin{tabular}{ll}
\toprule
Schema & Certificate\\
\midrule
Refl & \texttt{['r', t]}\\
Step & \texttt{['x', major, minor]}\\
Use$_*$ & \texttt{['u', index, antecedent\_proof]}\\
DE & \texttt{['e', goal, major, branch]}\\
DI & \texttt{['i', goal, witness, branch]}\\
\bottomrule
\end{tabular}
\end{center}
All array arities, tags, index types, term scopes, goal scopes, and rule
equalities are checked. In particular, Boolean values are not accepted
as integer indices. Use verifies the antecedent before opening the
clause body. Step checks the exact predecessor source and the successor
shape of its minor. Branch freshness is obtained by allocation, and
the branch result is checked against the target formula in the older
scope, preventing eigenparameter escape.

The JSON parser has its ordinary last-key behavior for duplicate textual
object keys; the contract is stated for the decoded object. As with the
comparison checkers, finite runtime, memory, and recursion limits may
reject sufficiently large valid finite proofs. The universal theorems
concern the corresponding finite-tree proof algorithms without such
resource bounds.

\subsection{Eliminating redundant scope state}

The closed standalone interface begins with an empty context and scope
zero. Each branch of DE or DI adds exactly one clause and allocates
exactly one eigenparameter. No other recursive call changes either.
Thus induction on the call tree gives
\[
          \texttt{scope}=\lvert\texttt{context}\rvert.
\]
The selected checker removes the separate scope argument and computes
it from the context. This is independent of deterministic specialization
and can also be applied to the unrestricted control. It is an invariant
of the closed public entry point; no equivalence is asserted for arbitrary
external calls to the older open-context helper with mismatched arguments.

\subsection{The executable proof transformation}

\path{code/translate.py} translates an entire certificate from the exact
incoming \path{reference/unrestricted.py} and verifies the result with
the primary checker. Its reverse mode inserts the selected target and
reverses the clause indices, then checks the expanded proof with the
unrestricted program. Both directions preserve the final formula
literally; the forward-then-reverse proof tree need not be the original
tree.

Clause occurrences are represented by distinct objects in the untrusted
compiler, so equal-looking clauses cannot be accidentally interchanged.
A mapping sends each source eigenlevel to the fresh target level chosen
under the present context embedding. A persistent premise is a function
that recursively recompiles the original proper subproof for the requested
extension. Formula transport consumes its D minor before entering new
branches. The auxiliary recursion follows the rank order in the proof;
it does not recursively translate the generated large proof.

The arithmetic constructors in \path{code/arith.py} are also untrusted.
They produce ordinary Refl/Step trees and never certify a theorem by
themselves. The checker imports neither file. This separates theorem
generation from the small acceptance boundary without omitting any
trusted checker component from the measurement.

\section{Controlled size measurements and verification}
\label{sec:measurements}

The environment is CPython 3.12.14. A true module docstring is removed,
the remaining complete module is rendered by \texttt{ast.unparse},
and one final LF is appended. Tokens exclude indentation, line-ending,
comment, encoding, and end-marker tokens. AST nodes count the parsed
module after docstring removal. Gzip uses level 9 and timestamp zero
and is reported for information, not as the optimization objective.
Raw source bytes and raw-source gzip sizes are also in
\path{evidence/metrics.json}; the primary readable source has 3597 bytes
and gzip size 1346 bytes, including its explanatory docstring.

\begin{center}
\begin{tabular}{lrrrr}
\toprule
Complete checker & Bytes & Tokens & AST & Gzip-9\\
\midrule
Unrestricted, explicit scope & 2950 & 829 & 902 & 1004 \\
Deterministic, explicit scope & 2938 & 829 & 903 & 999 \\
Unrestricted, scope from context & 2926 & 819 & 895 & 1000 \\
Deterministic, scope from context & 2914 & 819 & 896 & 997 \\
\bottomrule
\end{tabular}
\end{center}

The unrestricted matched controls store contexts newest first. This
changes no source-size coordinate relative to the original incoming
program; it changes gzip by one byte. The exact original source is
also included and measures 2950/829/902 with gzip 1005.

The deterministic rule itself removes 12 normalized bytes, leaves the
token count unchanged, and adds one AST node. The scope invariant removes
24 bytes, 10 tokens, and 7 AST nodes in either presentation. The combined
change therefore improves the explicit unrestricted control by 36 bytes,
10 tokens, and 6 AST nodes. The mathematical restriction should not be
described as a strict improvement in every coordinate by itself.

The historical compact 516-byte profile and the report's readable
checker have different public interfaces and validation boundaries.
Other filed typed presentations also use different computational
normalizers. We do not assert that 2914 is the repository-wide smallest
program. The separate 820-byte supplement addresses a more compact
unchanged native implementation on its own matched contract.

\subsection{Verification evidence}

The main suite contains 100 unrestricted positive certificates: 83
inherited uniform-Step examples and 17 independently constructed stress
cases. Every one was translated, accepted by the deterministic checker,
expanded back, and accepted with its identical goal by the original
checker. Each original, forward output, and reverse output also passed
the standalone command-line interface in ordinary and optimized Python:
600 complete process checks.

The independent operational audit contains 209 frozen certificates and
956 process checks. It compares explicit and implicit scope state,
checks 1200 proof mutations, rejects 93 altered false conclusions,
and independently evaluates 157 accepted closed equations in the natural
numbers. Its 25,047 observed explicit-scope calls all satisfy the scope
invariant. Dispatch and capture tests are designed to fail when the
chosen clause index or binder index is changed. The independent grammar
fusion suite has 22,350 public-entry shape comparisons and 248 process
checks. These numbers count executions, not disjoint mathematical
theorems or a formal proof of program correctness.

The compound examples exercise D ranks one, two, and three, numerical
targets containing older parameters, both coordinates bound by D,
parameters under iterator binders, and antecedent proofs that themselves
use compound clauses. These last cases test the persistence mechanism
that a simple weakening argument would miss.

\begin{center}
\small
\begin{tabular}{lrr}
\toprule
Example & Source JSON bytes & Target JSON bytes\\
\midrule
persistent\_compound\_rank\_3 & 699 & 175640 \\
rank3\_nested\_iterator & 1164 & 101884 \\
rank3\_older\_parameter & 1097 & 94509 \\
induction\_uses\_hypothesis & 324 & 1578 \\
\bottomrule
\end{tabular}
\end{center}

Certificate expansion is a substantial cost. A small checker need not
have short proofs or fast proof checking. No polynomial bound on the
translation's serialized output is claimed. The recursion proves
termination on each finite source proof; the tables report actual
outputs, not a general complexity upper bound.

\section{An independent 820-byte native checker refinement}
\label{sec:fusion}

This supplement preserves the original native scalar Eq/Q/D language
and its entire existing certificate format. Its exact 866-byte control
is \path{Research/scalar_uniqueness/code/compact/checker.py} at the frozen
commit, Git blob \texttt{00361c25eb276ab835103d18c767e80e85b6a04d}.
The full proof and portable audit are in \path{supplement/scope_fusion/}.

The old checker traverses a public goal once to validate its grammar and
again in the kernel's mandatory first scope traversal. The refinement
checks grammar during that first scope traversal. A flag is enabled only
at public entry and disabled before any proof packet is read.

At a list, the added condition is exactly the old arity/sort condition.
Its children receive the same depth adjustment. At an integer, the
added condition is the conjunction of the old parity/sort test and the
old scope test. Structural induction on the integer/list tree therefore
shows that the fused entry succeeds exactly when both old passes succeed;
on success it returns the same tree. The unchanged lexical input check
has already excluded other types, including in unused packet prefixes.
With the flag disabled, every added guard short-circuits and the remaining
kernel is the original one. Neither the tree nor the proof stream has
been mutated by the validation pass. Both programs consequently accept
the same finite certificates, apart from incidental finite resource
limits.

\begin{center}
\begin{tabular}{lrrrr}
\toprule
Complete native program & Bytes & Tokens & AST & Gzip-9\\
\midrule
Exact control & 866 & 381 & 439 & 504\\
Fused grammar/scope traversal & 820 & 360 & 412 & 493\\
\bottomrule
\end{tabular}
\end{center}
This is a 46-byte, 21-token, and 27-node improvement on that complete
contract. It changes no logical or arithmetic rule, so its unchanged
mathematical meaning is inherited directly. Its highly encoded source
is less readable than the primary presentation; the two programs serve
different comparison purposes. Neither one is advertised as a globally
minimal axiomatization.

\section{Reproduction, provenance, and limits}

From the extracted directory, run
\begin{lstlisting}
python3 reproduce.py
python3 code/checker.py < examples/induction.json
python3 code/translate.py < examples/compound_source.json
python3 reproduce.py --pdf
\end{lstlisting}
The first command runs the complete portable checks and refreshes the
source measurements. The last additionally needs a normal TeX Live
installation. The generated \path{report.tex} is self-contained:
it includes the full theorem, inherited appendices, and checker listing,
and can be compiled directly with pdfLaTeX. The readable assembly
fragments are retained under \path{paper/}.

\path{CLAIMS.json} distinguishes mathematical claims, executable evidence,
and limitations. \path{evidence/source_blobs.json} records the exact Git
blobs and SHA-256 hashes of the inherited bridge files; they match the
pinned repository byte for byte. \path{MANIFEST.json} hashes the final
distributed files. The original repository's MIT No Attribution license
is included. Newly written artifacts are supplied under the same
permissive terms for integration.

The new metatheorem is a written proof with independent mathematical
review and executable transformations. It has not been formalized in
Lean or Rocq. The repository's existing Lean developments support its
inherited core; their presence is not represented as verification of
our new deterministic specialization theorem. No new arithmetic-strength
claim, universal function evaluator, standard-number restriction,
quotient domain, or omitted logical rule is used in the argument.

\addtocontents{toc}{\protect\setcounter{tocdepth}{1}}
\appendix
\section{Inherited mathematical sources}
The next five appendices are the authenticated source fragments from
\path{Research/certified_zero/reference/} at the frozen commit. Their
historical references to ``new'' or ``current'' refer to those inherited
developments. They are included to provide the full logical and
arithmetical bridge, with attribution, and are not claimed as new work
of the present report.

% New mathematical result: complementary introduction choices.
% Requires amsmath/amssymb/amsthm and theorem, lemma, proposition environments.
% The inherited arithmetic-to-iPRA bridge must be supplied separately.

\section{Complementary introductions with one local assumption}
\label{sec:choice-introduction}

The old paired introduction retains both the zero instance and an
induction instance when its matrix is an equation. Both together are
unnecessary. A certificate may choose either one, provided the induction
choice is restricted to a literal equation. This changes neither the
formula language nor the meaning of a formula. The omitted combination
is recovered by a short nested introduction--elimination construction.

\subsection{Complete native syntax and rule inventory}

The scalar numerical language is
\[
 t::=x\mid0\mid St\mid I(n,a;u.h).
\]
Only $h$ binds the state variable $u$; $n$ and $a$ do not.
The update may contain fixed numerical parameters. The computation
axioms in every retained context are
\[
 I(0,a;u.h)=a,\qquad
 I(Sn,a;u.h)=h[I(n,a;u.h)/u].
 \tag{I0, IS}\label{eq:choice-computations}
\]
The argument also applies unchanged to the full fixed-description PR
language with its projection, composition, and recursion equations.
No numerical variable ranges over function descriptions, and no uniform
evaluator is postulated.

Native formulas and reusable context clauses are
\[
 A::=(s=t)\mid\mathsf D_{x,y}A,\qquad
 \langle A;y.B(y)\rangle,\quad y\notin\operatorname{FV}(A).
\]
Bound variables are renamed before capture-free substitution.
The two variables $x,y$ in a paired formula are distinct.
An ordinary assumption $F$ is the constant clause
$\langle H;v.F\rangle$, where $v$ is absent from $F$ and $H$ is the
fixed computation theorem $I(0,0;u.0)=0$. Thus ordinary hypothesis
use is a clause Use preceded by this computation proof.

There are computation axioms and exactly the following inference
patterns:
\[
\frac{\Gamma\vdash A}{\Gamma\vdash B(t)}
\quad\mathsf{Use}
\quad(\langle A;y.B(y)\rangle\in\Gamma),
\]
\[
\frac{\Gamma\vdash\mathsf D_{x,y}B(x,y)\qquad
      \Gamma,\langle B(a,0);y.B(a,y)\rangle\vdash F}
     {\Gamma\vdash F}
\quad\mathsf{DE},
\qquad
\frac{\Gamma\vdash s=t\qquad\Gamma\vdash F(s)}
     {\Gamma\vdash F(t)}
\quad\mathsf{Tr},
\]
\[
\frac{\Gamma,B(t,\eta_\varepsilon(z))\vdash B(t,Sz)}
     {\Gamma\vdash\mathsf D_{x,y}B(x,y)}
\quad\mathsf{DI}_{\varepsilon},
\qquad
\eta_0(z)=0,\quad\eta_1(z)=z.
\tag{Choice DI}\label{eq:choice-di}
\]
Here $\varepsilon$ is a certificate choice in $\{0,1\}$, and
$\varepsilon=1$ is permitted only when $B$ has a literal equation
root. An equation beneath a binder is not a literal equation.
Call $\varepsilon=0$ the \emph{case choice} and $\varepsilon=1$
the \emph{progress choice}. Both retain exactly one ordinary
formula assumption. Compound matrices have only the case choice.

In DI the eigenparameter $z$ is fresh for $\Gamma$, the witness
$t$, and every fixed parameter of $B$. In DE the eigenparameter
$a$ is fresh for $\Gamma$, the conclusion $F$, and every fixed
parameter of $B$. In transport, $F(v)$ is any fixed native
formula template. Displayed local assumptions and clauses are
discharged. There are no primitive reflexivity, coverage,
separation, ordinary logical, or additional induction rules.

We call this calculus $\mathcal C$. The original comparison
calculus $\mathcal N$ has identical syntax, computation,
Use, DE, and transport, but its introduction is
\[
\frac{\Gamma,B(t,0),[B(t,z)]_{\mathrm{at}}\vdash B(t,Sz)}
     {\Gamma\vdash\mathsf D_{x,y}B(x,y)}.
\tag{Old DI}\label{eq:choice-old-di}
\]
The bracket supplies $B(t,z)$ exactly at a literal equation.

\subsection{Exact equivalence and the unused-witness construction}

\begin{lemma}[Structural admissibility]
\label{lem:choice-structural}
Weakening, exchange, contraction, renaming, ordinary-assumption cut,
and capture-free numerical substitution are admissible in $\mathcal C$.
\end{lemma}
\begin{proof}
Induct on derivations. Rename local eigenparameters before inserting
proofs or terms. Numerical substitution preserves a formula's root
constructor, so an allowed progress choice remains allowed.
The selected choice does not change under substitution or cut.
For ordinary-assumption cut, replace each use of its constant clause
by the supplied proof. Computation schemas are closed under
capture-free numerical substitution. These facts verify every case.
\end{proof}

\begin{theorem}[Identity-language equivalence]
\label{thm:choice-equivalence}
For every retained clause context $\Gamma$ and native formula $F$,
\[
 \Gamma\vdash_{\mathcal C}F
 \quad\Longleftrightarrow\quad
 \Gamma\vdash_{\mathcal N}F.
\]
The transformations are effective on finite derivations and leave
all numerical and formula syntax unchanged.
\end{theorem}
\begin{proof}
Each new introduction is an instance of the old one after weakening
by the omitted local assumption. For the case choice at an equation,
add the induction instance. For the progress choice, add the zero
instance. Compound case introduction is already the old rule.
Thus every new derivation translates to an old one.

For the reverse direction, transform an old derivation recursively.
Only its atomic DI case needs a construction. Suppose its matrix is
the equation $E(x,y)$ and its transformed premise is
\[
 \pi:\quad\Gamma,E(t,0),E(t,z)\vdash E(t,Sz).
\]
The outer new inference will be case DI with witness $t$. It is
therefore enough to prove $E(t,Sz)$ under
$\Gamma,E(t,0)$.

Rename the eigenvariable in $\pi$ to a fresh $v$. Choose a new
variable $u$ absent from $E(t,y)$ and all retained data. Under
$\Gamma,E(t,0)$, progress DI with witness zero introduces
\[
             \mathsf D_{u,y}E(t,y)
\]
from exactly the renamed premise $\pi[v/z]$.
The bound coordinate $u$ is deliberately unused.

Apply DE to this paired formula. Because $u$ is unused, the opened
clause is exactly
\[
               \langle E(t,0);y.E(t,y)\rangle;
\]
it contains no new existential witness. The retained ordinary
assumption $E(t,0)$ supplies its antecedent. Use at $Sz$ gives
$E(t,Sz)$, and DE closes. Finally the outer case DI introduces
the original $\mathsf D_{x,y}E(x,y)$ with witness $t$.
The original premise proof $\pi$ has been inserted once.

An old compound introduction already is the new case choice.
All other rule cases are unchanged. This completes the induction
on the original derivation.
\end{proof}

The unused inner coordinate is essential. Introducing
$\mathsf D_{x,y}E(x,y)$ internally and then eliminating it would
give an arbitrary fresh witness, not necessarily the original
$t$. The construction instead keeps $t$ as a fixed parameter
of $E(t,y)$ and binds only the unused coordinate $u$.

\begin{proposition}[Linear bound on inference nodes]
\label{prop:choice-node-bound}
Count each computation axiom, Use, DI, DE, and transport as one
inference node, expanding ordinary-assumption use as a computation
of $H$ followed by Use. If an old derivation has $N$ nodes and
$r$ atomic DI nodes, the displayed transformation has exactly
\[
                   N+5r\leq 6N
\]
nodes.
\end{proposition}
\begin{proof}
Each atomic old DI node is replaced by six nodes: outer case DI,
inner progress DI, DE, Use of the opened clause, Use of the
retained base assumption, and the computation of $H$ needed for
that last ordinary-assumption use. Its old premise subtree
occurs once. All other inference nodes are copied once.
\end{proof}

This counts inference-tree nodes. It is not a bound on serialized
JSON bytes, formula sizes, checker runtime, or the number of
consumed packet cells when hypothesis closure is implicit.
In the formula-context Q/D presentation, the corresponding
replacement has outer case DI, inner progress DI, DE, QE,
and two Hyp nodes, giving the same six-node count.

\subsection{Induction and numerical cases separate cleanly}

\begin{lemma}[Contextual equation induction]
\label{lem:choice-ei}
The progress choice, DE, and Use derive
\[
\frac{\Gamma\vdash E(0)\qquad
      \Gamma,E(z)\vdash E(Sz)}
     {\Gamma\vdash E(t)}
\tag{EI}\label{eq:choice-ei}
\]
for every equation $E$, with the usual freshness condition.
\end{lemma}
\begin{proof}
Choose an unused witness coordinate $u$. Progress DI introduces
$\mathsf D_{u,y}E(y)$ from the given step. DE opens the clause
$\langle E(0);y.E(y)\rangle$. Supply its antecedent using the given
base proof and use the clause at $t$. Close DE.
\end{proof}

\begin{lemma}[Cases, equality, and binder congruence]
\label{lem:choice-preinduction}
The case choice derives, for arbitrary native formulas,
\[
\frac{\Gamma\vdash F(0)\qquad\Gamma\vdash F(Sz)}
     {\Gamma\vdash F(t)}
\tag{NC}\label{eq:choice-nc}
\]
with the usual freshness condition. Equality reflexivity,
symmetry, transitivity, numerical congruence, and uniform
matrix interderivability under $\mathsf D$ are derivable
without using the progress choice.
\end{lemma}
\begin{proof}
For cases, use case DI with an unused witness coordinate and
matrix $F(y)$, ignoring its retained base in the successor
premise. DE opens $\langle F(0);y.F(y)\rangle$; the given zero
proof and Use at $t$ finish.

For reflexivity put $r=I(0,t;u.0)$. Use the computation equation
$r=t$ as both transport premises with motive $v=t$ to obtain
$t=t$. The usual explicit equality motives give symmetry,
transitivity, and numerical congruence.

For binder congruence, open the old paired formula at a fresh
witness $a$ and introduce the new one by case DI at $a$.
Convert its retained base to the old base, use the opened
old clause at the successor, and convert back. Ordinary cut
closes the construction, and the reverse conversion gives
the converse. No induction is used.
\end{proof}

The primitive progress choice is still restricted to equations.
Allowing it at arbitrary compound matrices would, by the
argument for \eqref{eq:choice-ei}, supply induction for
arbitrary formulas and change the intended theory.

\subsection{The entire ordinary logical interface}
\label{subsec:choice-logic}

Here is a predecessor-free reconstruction of the ordinary logical
operations. It makes explicit that the case choice supplies the
logical interface independently of progress induction.

Let $C(v,a,b):=I(v,a;u.b)$ with $u$ unused. Computation gives
$C(0,a,b)=a$ and $C(Sv,a,b)=b$. Define the gate on fixed finite
formula syntax by
\[
 G_v(s=t):=(C(v,0,s)=C(v,0,t)),\qquad
 G_v(\mathsf D_{x,y}B):=\mathsf D_{x,y}G_vB,
\]
renaming bound variables away from the control. Equality,
binder congruence, and external structural induction give
\[
 \vdash G_0A,\qquad G_{Sv}A\dashv\vdash A,\qquad
 A\dashv\vdash\mathsf D_{u,y}G_yA
 \quad(u,y\notin\operatorname{FV}(A)).
\tag{Gate and padding}\label{eq:choice-gate}
\]
For padding, case DI proves the forward direction. DE, the
zero-gate theorem, and Use at $1$ prove the reverse direction.

Define a selector between two arbitrary formulas by
\[
 M_v(s=t,u=w):=(C(v,s,u)=C(v,t,w))
\]
and, with aligned fresh paired variables, by
\[
\begin{aligned}
 M_v(\mathsf D_{x,y}A,\mathsf D_{x,y}B)
   &:=\mathsf D_{x,y}M_v(A,B),\\
 M_v(E,\mathsf D_{x,y}B)
   &:=\mathsf D_{x,y}M_v(G_yE,B),\\
 M_v(\mathsf D_{x,y}A,E)
   &:=\mathsf D_{x,y}M_v(A,G_yE),
\end{aligned}
\]
where $E$ is an equation. Every call decreases maximum input
binder depth. Padding and congruence prove
$M_0(A,B)\dashv\vdash A$ and
$M_{Sv}(A,B)\dashv\vdash B$.

Put $\delta(v)=C(v,1,0)$ and $p(z)=I(z,0;u.\delta(u))$.
Only the fixed computations $p(1)=1$ and $p(2)=0$ are needed.
The full logical macro inventory is
\[
\begin{aligned}
 A\Rightarrow_*B&:=\mathsf D_{u,z}M_z(A,B),\\
 A\wedge_*B&:=\mathsf D_{u,z}G_zM_{p(z)}(A,B),\\
 \exists_*x\,A(x)&:=\mathsf D_{x,y}G_yA(x),\\
 A\vee_*B&:=\exists_*x\,M_x(A,B),\\
 \forall_*x\,A(x)&:=
        A(0)\wedge_*\mathsf D_{u,z}G_zA(z),\\
 \bot_*&:=(0=1).
\end{aligned}
\tag{Logical macros}\label{eq:choice-macros}
\]
Newly bound variables are fresh; $u$ is unused.

Implication introduction obtains its assumption from the zero
selector and supplies the consequent at successors, then uses
case DI. Elimination opens its clause and uses it at $1$.
Existential introduction uses its witness and the successor
gate; elimination opens at a fresh witness, supplies the
zero-gate theorem, and uses the clause at $1$.

Disjunction introduction uses witnesses $0$ and $1$.
For elimination, discharge the two branch proofs as
$A\Rightarrow_*F$ and $B\Rightarrow_*F$. NC applied to
$M_z(A,B)\Rightarrow_*F$ gives this implication at every
control: the two endpoint cases follow from the branch
implications. Existential elimination provides $M_a(A,B)$
and the implication gives $F$.

Given $A$ and $B$, selector endpoints and NC give $M_v(A,B)$
at an arbitrary control. Substitute $v=p(Sz)$ and use the
successor gate for conjunction introduction. For projections,
open its zero-gated clause and specialize at $1$ and $2$;
the two fixed computations give $B$ and $A$, respectively.

These rules prove existential numerical coverage:
\[
     (x=0)\vee_*\exists_*v\,(x=Sv).
\tag{Coverage}\label{eq:choice-coverage}
\]
Use NC. Zero takes the left injection of reflexivity, and a
successor takes the right injection with the evident witness.
Coverage gives equality cases even when the retained context
contains $x$; its elimination adds $x=0$ or $x=Sv$ rather
than silently changing that retained context.

Universal introduction follows from a fresh-parameter proof
by substituting zero and successors and forming the two
conjuncts. For elimination, take its $A(0)$ conjunct and
open the paired conjunct. Its base is a zero-gate theorem.
Coverage of the target chooses the zero instance from
$A(0)$ or a successor instance from the opened clause
and gate removal.

Finally, under $0=1$, computation and transport give
$s=C(0,s,t)=C(1,s,t)=t$ for arbitrary numerical terms.
Every equation follows, and external induction on formula
syntax with case DI gives every native formula. This
supplies falsity elimination. No excluded middle for
quantified formulas has entered the construction.

\subsection{Whole-domain realizations and both formula composites}

For a standard first-order reading, let $T_{\rm eq}$ have
intuitionistic equality logic, the same numerical signature
and computation equations, numerical coverage, separation
$(0=1)\to\bot$, and contextual equation induction.
Its interpretation of the native syntax is
\[
 (s=t)^\circ=(s=t),\qquad
 (\mathsf D_{x,y}B)^\circ
   =\exists x\bigl(B^\circ(x,0)\to\forall y B^\circ(x,y)\bigr),
\]
\[
 \langle A;y.B\rangle^\circ
       =A^\circ\to\forall y B^\circ.
\]
The domain is the entire numerical domain and interpreted
equality is actual equality.

Case DI is sound in this realization by numerical coverage.
Progress DI is sound by contextual equation induction under
the conditional base. Use, DE, transport, and computation
have their displayed ordinary meanings. Conversely, the
native macros derive every ordinary logical rule,
\eqref{eq:choice-ei} derives equation induction, and
\eqref{eq:choice-coverage} derives coverage. Separation
translates to the identity implication
$(0=1)\Rightarrow_*(0=1)$.

For the essential paired formula composite,
\[
 \mathsf D_{x,y}B
 \dashv\vdash
 \exists_*x\bigl(B(x,0)\Rightarrow_*\forall_*y\,B(x,y)\bigr).
\tag{Paired composite}\label{eq:choice-paired-composite}
\]
Forwards, DE opens its clause at a fresh $a$; under $B(a,0)$
that clause supplies $B(a,y)$ at a fresh arbitrary $y$.
Introduce the universal, implication, and existential
macros and close DE. Backwards, eliminate the existential
at a fresh $a$. Introduce the paired formula with witness
$a$ by case DI: the retained base, supplied implication,
and universal elimination at $Sz$ give its successor premise.

The ordinary interpretations of the six logical macros
have their ordinary round trips. For conjunction, its true
base and the instances $1,2$ recover both operands;
the converse uses selector cases. The universal has all
successor instances in its paired conjunct and zero in
its first conjunct; coverage gives exactly the ordinary
universal. Implication, disjunction, and existential
quantification use their selector or gate endpoints.
Falsity uses the explicit ordinary separation bridge.
External induction on finite formulas, together with
\eqref{eq:choice-paired-composite}, gives both formula
composites. Clauses correspond in both directions to
$A\Rightarrow_*\forall_*yB$ by the derived logical rules.

The numerical composites are literal identities with
internally provable graph $x=y$. No standard-numeral
restriction, retract, quotient, or equality reflection
assumption is being inserted at this stage. The separately
supplied arithmetic bridge identifies $T_{\rm eq}$ with
full iPRA; the present change preserves that bridge
unchanged because Theorem~\ref{thm:choice-equivalence}
preserves every native sequent.

\subsection{What the theorem does and does not claim}

The new inference interface retains one ordinary assumption
at every paired introduction. It preserves the entire
consequence relation, not just equations or closed
theorems. Consequently the established full-iPRA
interpretations and equiconsistency transfer unchanged.

This result is different from retaining only progress
introduction and strengthening closure elimination.
Here the original elimination rule is retained, and
the case and progress introductions coexist. It is also
different from allowing arbitrary compound progress,
which would supply stronger induction.

A complete-checker comparison must still count the
certificate choice, its equation-root restriction,
substitution, context handling, computation, and the
public input/output interface. The mathematical theorem
alone is not a byte-count or minimality theorem.
The inference-node bound is distinct from source-size
measurements and from serialized certificate size.
The proof is a written universal proof; finite regression
tests do not constitute a Lean/Rocq formalization of it.

\section{Arithmetic strength and both ordinary interpretation composites}
\label{sec:ordinary-bridge}

We spell out the arithmetic bridge to avoid identifying standard numerical
agreement with bi-interpretability. First elaborate every fixed iterator
into an ordinary function symbol. At stage zero the symbols are $0,S$.
Given a fixed earlier-stage term $h(u,\vec p)$, introduce
$i_h(n,a,\vec p)$ with the two equations
\[
i_h(0,a,\vec p)=a,\qquad
i_h(Sn,a,\vec p)=h(i_h(n,a,\vec p),\vec p).
\]
Only finite descriptions occur; the domain contains no code for a
universal evaluator. Formal parameters are distinct when a description
is declared, while actual arguments can coincide.

Let $T_I$ be ordinary intuitionistic first-order logic with actual
equality and genuine falsity, these staged equations, contextual
induction for literal equations, and
\[
\forall x\,(x=0\vee\exists y\,x=Sy),
\qquad (0=S0)\to\bot.
\]
The second sentence is the falsity bridge. It must be stated: positive
numerical equations by themselves admit a singleton structure. Native
explosion from $0=S0$ and a translation of ordinary falsity to that
equation do not authorize silently dropping the bridge from $T_I$.

The native-to-ordinary translation is
\[
(s=t)^\circ=(s=t),\qquad
(\mathsf D_{x,y}B)^\circ=
\exists x\bigl(B^\circ(x,0)\to\forall y B^\circ(x,y)\bigr).
\]
A clause $\langle A;y.B\rangle$ becomes $A^\circ\to\forall yB^\circ$.
For the auxiliary closure used by the compact checkers,
$(\mathsf Q_yB)^\circ=B^\circ(0)\to\forall yB^\circ$.
The interpretation domain is all elements, defined by $x=x$;
equality is actual equality.

Case DI is sound because numerical coverage combines its retained base
with its successor proof. Progress DI is sound by literal-equation
induction under a temporarily assumed base. Use, DE and transport are
ordinary logical inferences under the translations. Conversely the
derived logical operations, equation induction, cases and explosion in
Section~\ref{sec:choice-introduction} recover the rules and axioms of
$T_I$. An auxiliary closure is interderivable with
$\mathsf D_{v,y}B(y)$ for unused $v$: in one direction use case DI with
the closure and its base; in the other eliminate the paired formula
and return its supplied closure. Thus formulas and context clauses have
the stated meanings in both presentations.

\subsection{The arithmetic dependency order}

The scalar update has no recursion-counter argument. We do not conceal
one in an encoding. The inherited internal proofs, reproduced in the
appendices, establish the required arithmetic in this order:
\begin{enumerate}
\item Using equation induction and numerical cases, construct a fixed
scalar predecessor term and prove $P0=0$ and $P(Sn)=n$.
Existential coverage then gives $x=0\vee x=S(Px)$: in the successor
case $x=Sy$, apply $P$ to obtain $Px=y$ and substitute.
\item Construct truncated subtraction and a numerical characteristic
of actual equality. Prove its reflection on the entire original domain.
\item Compile each fixed quantifier-free formula into its characteristic
equation and recover every contextual quantifier-free induction instance.
\item With that induction available, construct pairing and projections
by a bounded square-root search and compile every fixed general
primitive-recursive description into a scalar iterator term.
\end{enumerate}
No pairing or faithful equality detector is assumed to obtain
predecessor. No appeal to standard-model correctness replaces an
internal equality. Every induction fixes its passive parameters;
diagonal substitution occurs only into a previously proved identity.

For a target recursion
\[
f(0,\vec p)=g(\vec p),\qquad
f(Sn,\vec p)=h(n,f(n,\vec p),\vec p),
\]
let $\pi,L,H$ be the proved pairing and projections. The compiler uses
\[
V(n,\vec p)=I\bigl(n,\pi(0,g^*(\vec p));
u.\pi(SL(u),h^*(L(u),H(u),\vec p))\bigr),
\qquad f^*(n,\vec p)=H(V(n,\vec p)).
\]
Equation induction proves $L(V(n,\vec p))=n$; the two projection
identities then prove the target defining equations. The construction
is external recursion on each finite PR description, not an internal
quantification over arbitrary functions.

\begin{theorem}[Whole-domain bi-interpretability and equiconsistency]
The ordinary realization $T_I$ and full first-order iPRA are
bi-interpretable. In both theories $x=y$ defines the isomorphism from
the original domain to its double interpretation. The native Choice-DI
calculus is therefore a faithful proof presentation of a theory
bi-interpretable and equiconsistent with iPRA.
\end{theorem}
\begin{proof}
Interpret an iterator in iPRA by primitive recursion that ignores its
counter argument in the update. iPRA proves its equations, literal
equation induction, coverage and separation. In the reverse direction
use the scalar compiler. The predecessor and equality-reflection proofs
supply the successor axioms and every quantifier-free induction instance.
All ordinary logical rules are preserved.

For two iterators with the same seed and provably equal updates,
equation induction on the equality of their values proves uniqueness at
every counter, with all parameters fixed. The corresponding argument
applies to ordinary primitive recursion. Induct externally on a finite
description: its double translation satisfies its original defining
equations, so internal uniqueness identifies the two functions. This
also proves coherence when substitution changes the staged syntactic
description of an update.

The domain and equality are unchanged; hence the identity graph is
internally total, bijective and preserves every interpreted operation.
The displayed ordinary meanings of the logical macros and the paired
binder give both formula composites by external induction on finite
formula syntax. Clauses are similarly equivalent to their conditional
universal meanings. The explicit falsity bridge preserves contradiction
in both directions. Effective translations of finite proofs give
equiconsistency, without a claim that either theory proves its own
consistency.
\end{proof}

% Adapted with attribution from the pinned unparenthesized one-hole archive.
\section{Internal scalar predecessor bootstrap}
\label{sec:hybrid-bootstrap}

This scalar bootstrap is inherited from the supplied scalar development~\cite{erasure} and is reproduced to keep the scalar checker self-contained. It is not a new contribution of this paper. The proof uses a factor map fixing odd states and a saturated-iteration construction of predecessor.

The available arithmetic terms are $0$, $S$, and
\[
 I(n,a;u.h),\qquad
 I(0,a;u.h)=a,\qquad
 I(Sn,a;u.h)=h[I(n,a;u.h)/u].
\]
The update binds only $u$.  All other variables in it are fixed numerical
parameters.  We use equality reasoning, the already reconstructed
zero/successor cases, and induction on a single literal equation.
Every induction below keeps its other parameters fixed.  Predecessor,
successor injectivity, pairing, and induction on arbitrary
quantifier-free formulas are not assumptions of the argument.

\subsection{Elementary terms and two finite cycles}

Write $1=S0$ and use larger numerals as successor abbreviations.  Define
\[
 C(n,a,b)=I(n,a;u.b),\qquad
 Z(n)=C(n,1,0),\qquad
 a+b=I(b,a;u.Su).
\]
The update variable of $C$ is absent from its body.  Consequently
\[
 C(0,a,b)=a,\qquad C(Sn,a,b)=b.
\]
Numerical cases show that $Z(n)$ is $0$ or $1$.
The usual commutative-monoid equations for addition follow by literal
equation induction.  For example, $0+b=b$ and $Sa+b=S(a+b)$ are
proved by induction on $b$; these yield commutativity by a further
equation induction, and associativity follows by induction on the last
summand.  We use these equations to reorder finite sums.  The notation
$2a$ means $a+a$ and $3a$ means $a+a+a$; no multiplication operation
is needed for this bootstrap.

Put
\[
 N(n)=I(n,0;u.Z(u)),\qquad
 j(u)=Z(u)+2N(u),\qquad
 M(n)=I(n,0;u.j(u)).
\]
The range of $N$ is contained in $\{0,1\}$: its zero value is $0$,
and its successor value is a value of $Z$.  This uses numerical cases
on its counter, not induction on a disjunction.  Hence
\begin{equation}
 N(n)+Z(N(n))=1,\qquad
 Z(Z(N(n)))=N(n),\qquad N(n+2)=N(n).
 \label{eq:hybrid-parity}
\end{equation}
Here the last equation follows from two successor computations and
the preceding equation.  Literal-equation inductions give
\begin{equation}
 N(2n)=0,\qquad N(2n+1)=1,\qquad N(3n)=N(n).
 \label{eq:hybrid-parity-scaling}
\end{equation}
For the first equation the successor step adds two to the argument of
$N$; the second follows by one computation.  In the third, increasing
$n$ adds three to the left argument and therefore applies $Z$ once to
both sides of the induction equation.

Finite computations now give
\[
 j(0)=1,\qquad j(1)=2,\qquad j(2)=0.
\]
The proof below needs just these three values and
$M(Sn)=j(M(n))$; it does not assume a division algorithm or a theorem
characterizing $M$ as remainder modulo three.

\subsection{A ramp with a global scaling identity}

Define
\[
 G(n)=I(n,0;u.u+M(u)+1).
\]
The central invariant is
\begin{equation}
 M(G(n))=N(n).
 \label{eq:hybrid-ramp-invariant}
\end{equation}
Its zero case is immediate.  For its successor step, retain the
single induction equation and split the Boolean value $N(n)$.
If it is $0$, the next $G$ state is $G(n)+1$; its $M$ value is
$j(0)=1$, equal to $N(Sn)$.  If it is $1$, the next $G$ state
is $G(n)+2$; two computations of $M$ change its value from $1$
to $2$ and then to $0$, again equal to $N(Sn)$.  This proves
the invariant by literal-equation induction.

It follows that
\[
 G(Sn)=G(n)+N(n)+1,\qquad G(n+2)=G(n)+3.
\]
Equation induction therefore yields
\begin{equation}
 G(2n)=3n.
 \label{eq:hybrid-ramp-double}
\end{equation}
A separate literal-equation induction proves
\begin{equation}
 G(3n)=3G(n)+N(n).
 \label{eq:hybrid-ramp-triple}
\end{equation}
For completeness, its successor calculation is
\begin{align*}
 G(3Sn)
 &=G(3n)+N(3n)+4
   =3G(n)+2N(n)+4,\\
 3G(Sn)+N(Sn)
 &=3G(n)+3N(n)+3+Z(N(n)).
\end{align*}
The expressions agree in both Boolean cases for $N(n)$.
The base is $0=0$.  In particular,
\eqref{eq:hybrid-ramp-triple} was proved internally; it was not
inferred from an intended floor-function interpretation of $G$.

Set
\[
 H(u)=C(N(u),G(u),u).
\]
The crucial choice is to retain $u$ when its parity is $1$.
The preceding equations give
\begin{equation}
 H(2u)=3u,\qquad H(1)=1,\qquad H(3u)=3H(u).
 \label{eq:hybrid-factor-map}
\end{equation}
To check the last identity, use $N(3u)=N(u)$ and split the Boolean
$N(u)$.  At $0$, both sides are $3G(u)$ by
\eqref{eq:hybrid-ramp-triple}; at $1$, both are $3u$.
This argument does not require that an arbitrary number has already
been represented as $2v$ or $2v+1$.

\subsection{A comparator from scalar iteration}

For every fixed update $h$, equation induction on $n$, keeping $m,a$
and all parameters fixed, proves the semigroup identity
\[
 I(m+n,a;h)=I(n,I(m,a;h);h).
\]
The base is the zero computation and the step applies $h$ to the
induction equation. This identity will also be used in the later
pairing construction.


For a fixed update $h$, equation induction on $y$, with its initializer
$a$ and all parameters of $h$ fixed, proves the rotation identity
\begin{equation}
 I(Sy,a;h)=I(y,h(a);h).
 \label{eq:hybrid-rotation}
\end{equation}
Both sides have base $h(a)$, and their successor equations apply the
same $h$ to the previous equality.  Here and below $I(y,a;h)$ is
an abbreviation for one fixed finite scalar term, not quantification
over functions.

By \eqref{eq:hybrid-factor-map}, another equation induction gives
\begin{equation}
 I(y,3a;H)=3I(y,a;H).
 \label{eq:hybrid-iterator-scaling}
\end{equation}
Its base is $3a=3a$.  In the step the fixed-parameter induction
equation rewrites the new left side as
\[
 H(3I(y,a;H))=3H(I(y,a;H)).
\]
The initializer parameter stays $a$ throughout this induction.
Equation induction also gives $I(y,1;H)=1$, since $H(1)=1$.

Define
\[
 E(x)=I(x,1;u.2u),\qquad
 L(x,y)=N\bigl(I(y,E(x);H)\bigr).
\]
Then the following entire recursion specification is provable:
\begin{equation}
 L(0,y)=1,\qquad L(Sx,0)=0,\qquad L(Sx,Sy)=L(x,y).
 \label{eq:hybrid-comparator}
\end{equation}
The first equation uses $E(0)=1$ and the fixed-point identity for
$H$.  The second uses $E(Sx)=2E(x)$ and $N(2a)=0$.
For the third,
\begin{align*}
 L(Sx,Sy)
 &=N\bigl(I(Sy,2E(x);H)\bigr)\\
 &=N\bigl(I(y,H(2E(x));H)\bigr)\\
 &=N\bigl(I(y,3E(x);H)\bigr)\\
 &=N\bigl(3I(y,E(x);H)\bigr)\\
 &=N\bigl(I(y,E(x);H)\bigr)=L(x,y).
\end{align*}
The successive steps use rotation, $H(2a)=3a$, iterator scaling,
and $N(3a)=N(a)$.  Thus $L$ has the comparison equations at all
arguments internally.  Its standard-model meaning as the indicator
of $x\le y$ is motivation, not a premise of the following argument.

\subsection{Saturated iteration gives predecessor}

We now need only the three equations
\eqref{eq:hybrid-comparator}.  Define
\[
 B_m(u)=C(L(Su,m),u,Su),\qquad
 W(t,m)=I(t,0;u.B_m(u)).
\]
The subscript in $B_m$ records a fixed parameter; $B_m$ is an open
term abbreviation, not a function variable.

Numerical cases on $e$ prove the selector equation
\[
 C(e,Sa,Sb)=S C(e,a,b).
\]
Equivalently one may use equation induction on $e$, whose induction
hypothesis is unused.  Therefore
\begin{align}
 B_{Sm}(Su)
 &=C(L(SSu,Sm),Su,SSu)\notag\\
 &=C(L(Su,m),Su,SSu)=S B_m(u),
 \label{eq:hybrid-conjugacy}\\
 B_{Sm}(0)&=C(L(S0,Sm),0,S0)=C(L(0,m),0,S0)=S0.
 \label{eq:hybrid-conjugacy-base}
\end{align}

Literal-equation induction on $t$, with $m$ fixed, proves
\begin{equation}
 W(St,Sm)=S W(t,m).
 \label{eq:hybrid-saturated-shift}
\end{equation}
The base is \eqref{eq:hybrid-conjugacy-base}.  The successor step is
\begin{align*}
 W(SSt,Sm)
 &=B_{Sm}(W(St,Sm))\\
 &=B_{Sm}(S W(t,m))\\
 &=S B_m(W(t,m))=S W(St,m).
\end{align*}
Only the second line uses the current induction hypothesis, at the
same fixed $m$.

Two subsequent equation inductions give
\begin{equation}
 L(Sn,n)=0,\qquad W(n,n)=n.
 \label{eq:hybrid-diagonals}
\end{equation}
The first uses the comparator recurrence.  The second uses
\eqref{eq:hybrid-saturated-shift} in its successor step.  The
substitution $t=m=n$ is made into that already proved universally
parameterized equation; no parameter of an active induction hypothesis
is changed.  Consequently
\begin{equation}
 W(Sn,n)=B_n(n)=C(L(Sn,n),n,Sn)=n.
 \label{eq:hybrid-extra-step}
\end{equation}

Finally define the finite scalar term
\begin{equation}
 P(m)=I\bigl(m,0;u.C(L(SSu,m),u,Su)\bigr).
 \label{eq:hybrid-predecessor-term}
\end{equation}
Its zero equation is immediate.  To justify its successor equation
without assuming a separate extensionality rule for iterator bodies,
put
\[
 V(t,m)=I\bigl(t,0;u.C(L(SSu,Sm),u,Su)\bigr).
\]
The comparator recurrence identifies this update pointwise with
$B_m(u)$.  Equation induction on $t$, with $m$ fixed, gives
$V(t,m)=W(t,m)$: their base values are $0$, and the successor step
uses the pointwise update equation and then the induction equation.
Now \eqref{eq:hybrid-extra-step} yields
\[
 \boxed{\qquad P(0)=0,\qquad P(Sm)=V(Sm,m)=W(Sm,m)=m.\qquad}
\]

\subsection{Dependency boundary}

All terms in this section are built from $0,S,I$ by finite substitution.
All inductions have a single equation as their invariant.  Numerical
cases and equality transport are the already reconstructed logical
resources.  Successor injectivity may now be obtained by applying $P$
to an assumed equation $Sx=Sy$.  Equality tests, reflection of
quantifier-free formulas, pairing, and the recovery of general
primitive recursion can then be developed in the established order.
They were not used in order to manufacture predecessor.

This argument is a written internal derivation in the stated theory.
The finite implementation checks do not constitute a proof-assistant
formalization of this entire metatheorem.

% Standalone section.  Requires only amsmath and amssymb in the parent document.
% No theorem environments, macros, or bibliography keys are assumed.
% All arithmetic claims in this section are proved here.

\section{From literal-equation induction to full intuitionistic PRA}
\label{sec:eic-completion}\label{app:arithmetic}

This section supplies an unconditional completion theorem. It is adapted
with attribution from the current forward-left draft~\cite{forward}. In the
scalar application used here, every function occurring in the proof is
already a scalar term: predecessor is supplied by
Appendix~\ref{sec:hybrid-bootstrap}, the selector is
$C(n,a,b)=I(n,a;u.b)$, addition is $x+y=I(y,x;u.Su)$, and
truncated subtraction is $d(x,y)=I(y,x;u.Pu)$. All other functions below
are finite compositions of these. Thus this proof does not presuppose
that arbitrary primitive recursion has already been compiled. The QFI
conclusion first applies to the scalar language; the general PR compiler
comes afterward.  Literal-equation
induction, together with the primitive-recursive defining equations,
predecessor coverage, and separation of $0$ and $1$, proves induction for every
quantifier-free formula.  In particular, the proof does not assume a faithful
numerical equality test, stability or decidability of equality, successor
uniqueness, or any commutation law for an iterator.  We construct the equality
test and prove its reflection property for arbitrary elements of the
numerical domain.

\subsection{The exact hypotheses}

Work in ordinary intuitionistic first-order logic with equality and genuine
falsity $\bot$.  The language has one numerical sort and a symbol for each
fixed, finite primitive-recursive description.  Distinct descriptions are
not identified merely because they compute the same standard function.
Let $\mathsf{PRE}$ contain the defining equations for zero, successor,
projections, composition, and primitive recursion.  In particular, choose
descriptions $P$ and $C$ with equations
\begin{align}
 P(0)&=0, & P(Sx)&=x, \label{eq:eic-predecessor}\\
 C(0,a,b)&=a, & C(Sx,a,b)&=b. \label{eq:eic-selector}
\end{align}
Write $1=S0$.  In addition to $\mathsf{PRE}$, assume
\begin{align}
 \mathsf{Cov}:&\qquad
       \forall x\,\bigl(x=0\lor x=S(Px)\bigr),
       \label{eq:eic-coverage}\\
 \mathsf{Sep}:&\qquad 0=1\longrightarrow\bot,
       \label{eq:eic-separation}
\end{align}
and induction for each \emph{literal equation} $E(n,\vec p)$:
\begin{equation}
 \forall\vec p\,
 \left[
   \left(E(0,\vec p)\land
       \forall n\bigl(E(n,\vec p)\longrightarrow E(Sn,\vec p)\bigr)
   \right)
   \longrightarrow\forall n\,E(n,\vec p)
 \right].
 \label{eq:eic-induction}
\end{equation}
Here the schematic expression $E$ is a single equality of numerical
terms; the displayed universal quantifiers range only over numbers.
Every description and every instance of the schema is finite external
syntax.

The axiom formulation \eqref{eq:eic-induction} and the usual contextual
induction rule are interchangeable.  For example, if
\[
 \Gamma\vdash E(0,\vec p),
 \qquad
 \Gamma,E(z,\vec p)\vdash E(Sz,\vec p),
\]
where $z$ is absent from $\Gamma$ and the passive parameters, discharge the
second premise's equation and universally introduce $z$.  The axiom then
gives $\Gamma\vdash\forall n E(n,\vec p)$, which can be instantiated at any
target term.  Conversely, assume the antecedent of the axiom and apply the
contextual rule with fresh variables.  This uses the ordinary intuitionistic
logical rules and permits arbitrary, including quantified, formulas in
$\Gamma$.

\subsection{Constructor separation and two elementary tests}

Define the sign description and the countdown description by
\begin{equation}
 \operatorname{sg}(u)=C(u,0,1),
 \qquad
 K(u)=C(u,1,Pu).
 \label{eq:eic-sign-countdown}
\end{equation}
Their defining equations imply
\begin{align}
 \operatorname{sg}(0)&=0,
 &\operatorname{sg}(Sx)&=1,\label{eq:eic-sign-computation}\\
 K(0)&=1,
 &K(Sx)&=x.\label{eq:eic-countdown-computation}
\end{align}

\paragraph{Constructor laws.}
Successor is injective: applying $P$ to $Sx=Sy$ gives $x=y$ by
\eqref{eq:eic-predecessor}.  Successor never equals zero.  Indeed,
$Sx=0$ gives $1=0$ by applying $\operatorname{sg}$ and using
\eqref{eq:eic-sign-computation}; symmetry and
\eqref{eq:eic-separation} give $\bot$.  Consequently $1\ne0$, where
$A\ne B$ abbreviates $A=B\to\bot$.  None of these arguments uses
induction.

\paragraph{The sign test.}
For every $u$,
\begin{align}
 &\operatorname{sg}(u)=0\lor\operatorname{sg}(u)=1,
       \label{eq:eic-sign-range}\\
 &\operatorname{sg}(u)=0\quad\longleftrightarrow\quad u=0,
       \label{eq:eic-sign-zero}\\
 &\operatorname{sg}(u)=1\quad\longleftrightarrow\quad u\ne0.
       \label{eq:eic-sign-one}
\end{align}
To prove the range assertion, split \eqref{eq:eic-coverage}.  The branch
$u=0$ gives sign $0$, and the branch $u=S(Pu)$ gives sign $1$.
The implication $u=0\to\operatorname{sg}(u)=0$ is computation.
For its converse, the successor branch of coverage would make
$\operatorname{sg}(u)=0$ imply $1=0$, a contradiction; in that branch
explosion gives $u=0$ as well.  Similarly, if
$\operatorname{sg}(u)=1$ and $u=0$, computation gives $0=1$.
Conversely, $u\ne0$ excludes the zero branch of coverage, and the
successor branch computes the sign as $1$.

Coverage and \eqref{eq:eic-selector} also give, for all $u,v$,
\begin{equation}
 C(u,v,v)=v.
 \label{eq:eic-selector-idempotence}
\end{equation}
Both coverage branches compute to $v$, so this identity again needs no
induction.

\paragraph{The countdown test and the notation $W$.}
For every $u$,
\begin{equation}
 K(u)=0\quad\longleftrightarrow\quad u=1.
 \label{eq:eic-countdown-zero}
\end{equation}
If $u=1=S0$, \eqref{eq:eic-countdown-computation} gives $K(u)=0$.
Conversely, assume $K(u)=0$ and split coverage.  If $u=0$, then
$K(u)=1$, contradicting separation, and explosion gives $u=1$.
If $u=S(Pu)$, computation gives $K(u)=Pu$; hence $Pu=0$ and
$u=S0=1$.

If a presentation uses a description named $W$ with defining equations
\[
 W(0)=1,\qquad W(Sx)=x,
\]
then it proves $K(u)=W(u)$ for every $u$ already from coverage.
In the zero branch both sides are $1$; in the branch $u=S(Pu)$ both
sides are $Pu$.  Thus $K$ and $W$ are provably equal on the whole domain,
even if they are distinct descriptions.  This equality does not appeal to
standard-model agreement or to uniqueness of an arbitrary recurrence.

\subsection{Four elementary uses of equation induction}

Choose descriptions for addition and truncated subtraction, with recursion
on the second argument:
\begin{align}
 x+0&=x, &x+Sy&=S(x+y),\label{eq:eic-addition}\\
 d(x,0)&=x, &d(x,Sy)&=P(d(x,y)).\label{eq:eic-subtraction}
\end{align}
Using these equations and \eqref{eq:eic-induction}, we now derive
\begin{align}
 0+y&=y,\label{eq:eic-left-zero}\\
 (Sx)+y&=S(x+y),\label{eq:eic-left-successor}\\
 d(Sx,Sy)&=d(x,y),\label{eq:eic-cancel-successors}\\
 d(Sa,a)&=1.\label{eq:eic-diagonal}
\end{align}
Each induction below has a single equation as its induction formula.

For \eqref{eq:eic-left-zero}, induct on $y$.  The base is $0+0=0$.
The step is $0+Sy=S(0+y)=Sy$.

For \eqref{eq:eic-left-successor}, hold $x$ fixed and induct on $y$.
At zero both sides are $Sx$.  At a successor the induction hypothesis
gives
\[
 (Sx)+Sy=S\bigl((Sx)+y\bigr)
       =S\bigl(S(x+y)\bigr)=S(x+Sy).
\]
Combining \eqref{eq:eic-left-successor} with the defining equation for
addition gives the useful rearrangement
\begin{equation}
 (Sy)+r=y+Sr.
 \label{eq:eic-shift-successor}
\end{equation}
Indeed, both sides are $S(y+r)$.

For \eqref{eq:eic-cancel-successors}, hold $x$ fixed and induct on $y$.
The base computes as
\[
 d(Sx,S0)=P(d(Sx,0))=P(Sx)=x=d(x,0).
\]
Applying $P$ to the induction equation gives the step:
\[
 d(Sx,S(Sy))=P(d(Sx,Sy))=P(d(x,y))=d(x,Sy).
\]

Finally, induct on $a$ for \eqref{eq:eic-diagonal}.  The base is
$d(S0,0)=S0=1$.  At a successor,
\[
 d(S(Sa),Sa)=d(Sa,a)=1,
\]
using \eqref{eq:eic-cancel-successors} and the induction hypothesis.

\subsection{An equation invariant that reflects actual equality}

The remaining arithmetic ingredient is the following equation, valid for
all $x,y$:
\begin{equation}
 C\bigl(\operatorname{sg}(d(x,y)),\ x,\ y+d(x,y)\bigr)=x.
 \label{eq:eic-residual-invariant}
\end{equation}
Its meaning is that a nonzero subtraction residual reconstructs the
original argument by addition.  The selector expresses this assertion as
one literal equation, so equation induction applies directly.

Hold $x$ fixed and induct on $y$.  At $y=0$, subtraction gives residual
$x$, and \eqref{eq:eic-left-zero} gives $0+x=x$.  The desired equation
is consequently $C(\operatorname{sg}(x),x,x)=x$, which is
\eqref{eq:eic-selector-idempotence}.

For the step, assume \eqref{eq:eic-residual-invariant} at $y$ and
abbreviate
\[
 r=d(x,y),\qquad r'=Pr=d(x,Sy).
\]
These are local term abbreviations, not additional function symbols.
Split coverage for $r'$.

If $r'=0$, then $\operatorname{sg}(r')=0$, and the successor instance
of \eqref{eq:eic-residual-invariant} reduces immediately to $x=x$.

In the other branch, $r'=S(Pr')$, so constructor separation gives
$r'\ne0$.  Split coverage for $r$.  Its zero branch would imply
$r'=P0=0$, a contradiction, from which the desired step equation follows.
In its successor branch,
\[
 r=S(Pr)=Sr',\qquad \operatorname{sg}(r)=1.
\]
The induction equation therefore computes to $y+r=x$.  Meanwhile,
\eqref{eq:eic-shift-successor} gives
\[
 (Sy)+r'=y+Sr'=y+r=x.
\]
Since $r'$ is a successor, $\operatorname{sg}(r')=1$, so this is
exactly the branch selected in the successor instance of
\eqref{eq:eic-residual-invariant}.  This completes the equation induction.
All case distinctions came from coverage; no induction on an implication
or a disjunction was used.

\paragraph{Reflection at residual $1$.}
We have, for arbitrary $a,b$,
\begin{equation}
 d(Sa,b)=1\quad\longleftrightarrow\quad a=b.
 \label{eq:eic-residual-reflection}
\end{equation}
For the forward direction, substitute $x=Sa$, $y=b$ into
\eqref{eq:eic-residual-invariant}.  If $d(Sa,b)=1$, the sign is $1$
and the selector returns its second branch, giving
\[
 b+1=Sa.
\]
The defining equation for addition gives $b+1=S(b+0)=Sb$.  Thus
$Sb=Sa$, and applying $P$ gives $b=a$, hence $a=b$.
Conversely, $a=b$ and \eqref{eq:eic-diagonal} give
$d(Sa,b)=d(Sa,a)=1$.

Define the equality characteristic by the fixed description
\begin{equation}
 e(a,b)=\operatorname{sg}\bigl(K(d(Sa,b))\bigr).
 \label{eq:eic-equality-characteristic}
\end{equation}
The sign, countdown, and residual reflection lemmas give
\begin{align}
 e(a,b)=0&\quad\longleftrightarrow\quad a=b,
       \label{eq:eic-equality-reflection}\\
 e(a,b)=0&\ \lor\ e(a,b)=1.
       \label{eq:eic-equality-range}
\end{align}
In detail, $e(a,b)=0$ is equivalent to $K(d(Sa,b))=0$ by
\eqref{eq:eic-sign-zero}, which is equivalent to $d(Sa,b)=1$ by
\eqref{eq:eic-countdown-zero}, and then to $a=b$ by
\eqref{eq:eic-residual-reflection}.  The range assertion is
\eqref{eq:eic-sign-range} applied to $K(d(Sa,b))$.

The reflection property concerns the logic's actual equality throughout
the domain.  It was not assumed as part of a numerical comparison symbol,
and its proof did not restrict $a$ or $b$ to standard numerals.
It also yields decidable equality.  In the $e(a,b)=0$ branch,
\eqref{eq:eic-equality-reflection} gives $a=b$.  In the $e(a,b)=1$
branch, an assumption $a=b$ would give $e(a,b)=0$ and therefore $0=1$,
contradicting separation.  Hence
\begin{equation}
 \forall a\,\forall b\,\bigl(a=b\lor(a=b\to\bot)\bigr)
 \label{eq:eic-equality-decidable}
\end{equation}
is a theorem, not an extra hypothesis.

\subsection{Compilation of all quantifier-free formulas}

For each fixed quantifier-free formula $A(\vec p)$, construct a term
$b_A(\vec p)$ by external recursion on its finite syntax.  Take equality,
$\bot$, $\land$, $\lor$, and $\to$ as the primitive formula constructors;
negation is implication into $\bot$.  Set
\begin{align}
 b_{s=t}&=e(s,t), & b_{\bot}&=1,\label{eq:eic-boolean-atoms}\\
 b_{A\land D}&=C(b_A,b_D,1),\label{eq:eic-boolean-and}\\
 b_{A\lor D}&=C(b_A,0,b_D),\label{eq:eic-boolean-or}\\
 b_{A\to D}&=C(b_A,b_D,0).\label{eq:eic-boolean-implies}
\end{align}
Every term here is a finite composition of fixed primitive-recursive
descriptions; if desired, projections and composition give it a single
description symbol.  No internal quantification over formulas or programs
is involved.

An external induction on the syntax of $A$ proves, within the arithmetic
theory,
\begin{equation}
 \bigl(b_A=0\lor b_A=1\bigr)
 \quad\text{and}\quad
 \bigl(b_A=0\longleftrightarrow A\bigr).
 \label{eq:eic-boolean-correctness}
\end{equation}
Here are the full connective checks.  In every case $b_A=1$ implies
$\neg A$: if $A$ held, the already established implication
$A\to b_A=0$ would contradict separation.

\paragraph{Atomic formulas and falsity.}
For an equation, both claims are
\eqref{eq:eic-equality-reflection}--\eqref{eq:eic-equality-range}.
For falsity the range assertion is $1=0\lor1=1$, with its right disjunct
proved by reflexivity.  Also $1=0\to\bot$ follows from separation and
symmetry, while $\bot\to1=0$ follows by explosion.  Thus $1=0$ is
equivalent to $\bot$, as required.

\paragraph{Conjunction.}
Split $b_A=0\lor b_A=1$.  In the zero branch, $A$ holds and
$b_{A\land D}=b_D$.  Its range follows from the range of $b_D$,
and its being zero is equivalent to $D$, hence to $A\land D$ under
the known $A$.  In the one branch, $\neg A$ holds and
$b_{A\land D}=1$.  Its range is immediate.  If it were zero,
separation would give $\bot$, hence $A\land D$; conversely an
assumption $A\land D$ contradicts $\neg A$, and explosion gives
$b_{A\land D}=0$.

\paragraph{Disjunction.}
In the branch $b_A=0$, the formula $A$ holds and
$b_{A\lor D}=0$.  The range is immediate, and both $b_{A\lor D}=0$
and $A\lor D$ hold, proving their equivalence.  In the branch
$b_A=1$, we have $\neg A$ and $b_{A\lor D}=b_D$.  Its range follows
from the range of $b_D$.  If its value is zero, then $D$ and hence
$A\lor D$.  Conversely, from $A\lor D$, the $A$ alternative
contradicts $\neg A$ and yields $b_D=0$ by explosion, while the
$D$ alternative yields $b_D=0$ directly.  This proves the required
equivalence in both alternatives.

\paragraph{Implication.}
In the branch $b_A=0$, the formula $A$ holds and
$b_{A\to D}=b_D$.  Its range follows from the range of $b_D$.
If its value is zero, then $D$, which gives $A\to D$.  Conversely,
$A\to D$ and the known $A$ give $D$, hence $b_D=0$.
In the branch $b_A=1$, we have $\neg A$ and
$b_{A\to D}=0$.  Its range is immediate.  To prove $A\to D$,
assume $A$; this contradicts $\neg A$, and explosion gives $D$.
Thus both $b_{A\to D}=0$ and $A\to D$ hold in this branch, and
their equivalence follows.  This argument uses the proved two-value
range of the finite characteristic.  It does not use excluded middle
for arbitrary formulas or change the surrounding intuitionistic logic.

This completes the external syntactic induction proving
\eqref{eq:eic-boolean-correctness} for every quantifier-free formula.

\subsection{Recovery of induction and the full theory}

Let $A(n,\vec p)$ be any fixed quantifier-free formula.  Suppose an
arbitrary retained context $\Gamma$ proves its base and its induction
step:
\[
 \Gamma\vdash A(0,\vec p),
 \qquad
 \Gamma,A(z,\vec p)\vdash A(Sz,\vec p),
\]
with the usual freshness condition on $z$.  By
\eqref{eq:eic-boolean-correctness}, the base gives
$\Gamma\vdash b_A(0,\vec p)=0$.  Under the additional assumption
$b_A(z,\vec p)=0$, the same equivalence gives $A(z,\vec p)$,
the given step gives $A(Sz,\vec p)$, and the equivalence gives
$b_A(Sz,\vec p)=0$.  Apply literal-equation induction to
$b_A(n,\vec p)=0$ and then apply its reflection implication.
The result is $\Gamma\vdash A(t,\vec p)$ for every target term $t$.
Thus the theory proves the full contextual quantifier-free induction
rule, including in arbitrary quantified contexts.

\paragraph{Completion theorem.}
In the common language of fixed primitive-recursive descriptions, the
theory
\[
 \mathsf{PRE}+\mathsf{Cov}+\mathsf{Sep}+\mathsf{EI}
\]
has exactly the same sequents as ordinary first-order intuitionistic
primitive recursive arithmetic with all quantifier-free induction
instances.  The forward simulation follows from the recovered
quantifier-free induction, successor injectivity, and successor separation;
all primitive-recursive defining equations and all logical rules are
shared.  For the reverse simulation, literal-equation induction is an
instance of quantifier-free induction, and successor separation gives
$\mathsf{Sep}$.  Quantifier-free induction proves $\mathsf{Cov}$:
at zero choose its left disjunct, and at a successor use
$Sx=S(P(Sx))$, which follows from $P(Sx)=x$, for its right disjunct.

Both interpretations therefore use the full domain, defined by $x=x$,
and preserve every function symbol, equality, logical connective, and
quantifier.  Both composites have the identity isomorphism, with graph
$x=y$.  The theories are consequently identity-domain bi-interpretable
and equiconsistent.  This completion theorem is unconditional once the
literal-equation induction rule has been obtained; it has no dependency
on a classical successor-uniqueness reduction.

\section{Pairing, primitive recursion, and internal numerical composites}
\label{sec:iterator-coding}
This appendix is adapted with attribution from the progress-closure
arithmetic development~\cite{progress}. Its prerequisites have now been
proved: scalar predecessor, arithmetic reflection, and full quantifier-free
induction. Pairing is a derived proof tool; it is not part of any checker.

\subsection{Pairing from a bounded square-root search}

Once QFI is available, a familiar square-based injection gives a
particularly direct compiler. Write $a\le b$ for $d(a,b)=0$.
The usual linear-order and semiring identities follow by fixed
quantifier-free inductions from the successor, addition, and monus
equations. Multiplication itself is the scalar term
\[
                        a\cdot b=I(b,0;u.u+a).
\]
Only primitive-recursive arithmetic is involved in the following
bounded root search; there is no minimization over an unbounded domain.

For a fixed parameter $n$, put
\[
 U_n(v)=C(d((Sv)^2,n),Sv,v),\qquad
 R_n(i)=I(i,0;v.U_n(v)),\qquad q(n)=R_n(n).
\]
The update increments $v$ exactly when $(v+1)^2\le n$, and
otherwise retains it. QFI in $i$, with $n$ fixed, proves
\[
 R_n(i)\le i,\qquad R_n(i)^2\le n,\qquad
 R_n(i)=i\ \vee\ n<(S(R_n(i)))^2.              \tag{Root invariant}
\]
The base has $R_n(0)=0$. If the update increments, its test
contradicts the strict-upper-bound alternative in the third invariant.
Thus the old state equals $i$, and the new state equals $Si$;
its square bound is the increment test. If the update does not increment, the
strict upper square bound holds and the state remains fixed.
An already established strict upper square bound forces every
subsequent update to remain fixed. These are exactly the cases
needed for the three displayed invariants.

At $i=n$, either the strict upper square bound already holds or
$R_n(n)=n$, in which case $n<(n+1)^2$ supplies it. Thus
\[
                        q(n)^2\le n<(q(n)+1)^2.           \tag{Root}
\]
Monotonicity of squares shows uniqueness of a number satisfying
Root. All formulas just used are quantifier free in the scalar
language, with fixed $n$ during the induction.

Now define
\[
 \pi(a,b)=(a+b)^2+a,\qquad
 L(n)=d(n,q(n)^2),\qquad H(n)=d(q(n),L(n)).
\]
For $k=a+b$, one has $k^2\le k^2+a<(k+1)^2$ since $a\le k$.
Root uniqueness gives $q(\pi(a,b))=a+b$, and therefore
\[
                    L(\pi(a,b))=a,\qquad H(\pi(a,b))=b.
\]
Surjectivity of $\pi$ is unnecessary. The original numerical domain
is never replaced by the range of this coding; pairing is used only
to store the finite state of a fixed primitive recursion. This
standard square pairing is a convenient reconstruction here, not
a claimed new arithmetic basis.

\subsection{Compiling each primitive recursion}

Suppose the fixed target description is
\[
 f(0,\vec p)=g(\vec p),\qquad
 f(Sn,\vec p)=h(n,f(n,\vec p),\vec p).
\]
Assume by external structural induction that the earlier $g,h$
already have scalar translations $g^*,h^*$. Define
\[
 W(u,\vec p)=\pi(SL(u),h^*(L(u),H(u),\vec p)),
\]
\[
 V(n,\vec p)=I(n,\pi(0,g^*(\vec p));u.W(u,\vec p)),
 \qquad f^*(n,\vec p)=H(V(n,\vec p)).
\]
The update has only its state variable bound, with every $p_i$
fixed. Equation induction proves $L(V(n,\vec p))=n$.
The two projection equations and scalar computation then prove
both defining equations for $f^*$ internally. Initial functions
and composition translate directly. This completes the translation
of every fixed finite PR description without adding a universal
evaluator.

\subsection{Why both numerical round trips are internal}

Two fixed scalar iterators with the same seed and provably equal
updates are internally equal at all counters: use literal-equation
induction comparing their values. Their base agrees, and their
successor equations plus congruence prove the step. The same
uniqueness argument works for ordinary primitive recursion.

For an ordinary PR description, induct externally on its finite
construction. The double translation satisfies its original
defining equations, so internal equation induction identifies it
with the original operation. For a native scalar description,
earlier translated updates agree by the external induction
hypothesis; scalar uniqueness identifies the double-translated
iterator. These proofs hold at arbitrary internal counters.

Thus the domain relation $x=y$ is provably a total bijection
preserving every fixed numerical operation in both composites.
The formula constructions and round trips in
Section~\ref{sec:choice-introduction} give the corresponding formula
composites. Both interpretations preserve falsity, taking native
$0=1$ to an ordinary contradiction and ordinary $\bot$ to native
$0=1$. Finite proof translations therefore establish equiconsistency.


\clearpage
\section{Complete primary checker}
The following is the entire primary acceptance program after the stated
normalization. The explanatory module docstring is omitted from this
listing and the size comparison; it is retained in the readable source.
\begin{lstlisting}[language=Python]
import json
import sys

def require(condition):
    if not condition:
        raise ValueError('Invalid certificate')

def iterator(n, a, h):
    return a if n == 0 else ['i', n, a, h]

def term(a, scope, depth=0):
    if type(a) is int:
        return -scope <= a <= depth
    if type(a) is not list:
        return False
    if len(a) == 2 and a[0] == 's':
        return term(a[1], scope, depth)
    return len(a) == 4 and a[0] == 'i' and (a[1] != 0) and term(a[1], scope, depth) and term(a[2], scope, depth) and term(a[3], scope, depth + 1)

def formula(a, scope, depth=0):
    return type(a) is list and (len(a) == 3 and a[0] == '=' and all((term(t, scope, depth) for t in a[1:])) or (len(a) == 2 and a[0] == 'd' and formula(a[1], scope, depth + 2)))

def opening(a, value, depth=0):
    if type(a) is int:
        return value if a == depth + 1 else a - (a > depth + 1)
    tag = a[0]
    b = [opening(t, value, depth + (2 if tag == 'd' else int(tag == 'i' and i == 2))) for i, t in enumerate(a[1:])]
    return iterator(*b) if tag == 'i' else [tag] + b

def derive(proof, context):
    scope = len(context)
    match proof:
        case ['u', index, premise]:
            require(type(index) is int and 0 <= index < len(context))
            antecedent, body = context[index]
            require(derive(premise, context) == antecedent)
            return opening(body, -scope if index else ['s', 0])
        case ['r', a]:
            require(term(a, scope))
            return ['=', a, a]
        case ['x', major, minor]:
            match (derive(major, context), derive(minor, context)):
                case [['=', s, t], ['=', ['i', ['s', n], a, h], r]]:
                    require(s == iterator(n, a, h))
                    return ['=', r, opening(h, t)]
        case ['e' | 'i' as tag, goal, data, branch]:
            require(formula(goal, scope))
            fresh = -scope - 1
            if tag == 'e':
                paired = derive(data, context)
                require(paired[0] == 'd')
                body = opening(paired[1], fresh, 1)
                clause = (opening(body, 0), body)
                target = goal
            else:
                require(goal[0] == 'd' and term(data, scope))
                body = opening(goal[1], data, 1)
                eta = fresh if goal[1][0] == '=' and data == 0 else 0
                clause = (['=', 0, 0], opening(body, eta))
                target = opening(body, ['s', fresh])
            require(derive(branch, [clause] + context) == target)
            return goal
    raise ValueError('Invalid proof')

def verify(data):
    require(type(data) is dict and set(data) == {'goal', 'proof'})
    require(formula(data['goal'], 0))
    require(derive(data['proof'], []) == data['goal'])
    return True
if __name__ == '__main__':
    try:
        verify(json.load(sys.stdin))
    except Exception:
        sys.exit(1)
    print('ACCEPT')
\end{lstlisting}

\clearpage
\begin{thebibliography}{99}
\raggedright
\bibitem{baseline}
Vladimir Reshetnikov, \emph{ProveIt: Concise presentations of
intuitionistic primitive recursive arithmetic}, commit
\texttt{5d85023b0688e03b75f0ec3a04bb43ef0bb51705}.
\href{https://github.com/VladimirReshetnikov/ProveIt/tree/5d85023b0688e03b75f0ec3a04bb43ef0bb51705/Logic/PrimitiveRecursiveArithmetic/ConciseIPRA}{Frozen ConciseIPRA subtree}.
\bibitem{uniform}
ProveIt incoming research package, \emph{One successor step for equality
and computation}, \path{docs/incoming/ipra_uniform_step.zip}, at the
commit in~\cite{baseline}. Exact source checker and inherited test
certificates included in this package.
\bibitem{rootfree}
ProveIt incoming research package, \emph{A smaller clausal presentation
of iPRA}, \path{docs/incoming/ipra_rootfree_clausal_2026-11-11.zip}, at
the commit in~\cite{baseline}.
\bibitem{reduced}
ProveIt incoming research package, \emph{Reduced clause calculus},
\path{docs/incoming/ipra_reduced_clause_calculus_2026-11-11.zip}, at the
commit in~\cite{baseline}.
\bibitem{choice}
ProveIt research package, \emph{Complementary introductions with one
local assumption} and \emph{Arithmetic strength and both ordinary
interpretation composites},
\path{Research/certified_zero/reference/inherited_choice_basis.tex}
and \path{inherited_ordinary_bridge.tex}, at the commit in~\cite{baseline}.
Reproduced with attribution.
\bibitem{erasure}
ProveIt research package, \emph{Internal scalar predecessor bootstrap},
\path{Research/certified_zero/reference/scalar_predecessor.tex}, at the
commit in~\cite{baseline}. Reproduced with attribution.
\bibitem{forward}
ProveIt research package, \emph{From literal-equation induction to full
intuitionistic PRA},
\path{Research/certified_zero/reference/equality_reflection.tex}, at
the commit in~\cite{baseline}. Reproduced with attribution.
\bibitem{progress}
ProveIt research package, \emph{Pairing, primitive recursion, and internal
numerical composites},
\path{Research/certified_zero/reference/primitive_recursion.tex}, at
the commit in~\cite{baseline}. Reproduced with attribution.
\bibitem{compact}
ProveIt research package, \emph{Scalar uniqueness},
\path{Research/scalar_uniqueness/code/compact/checker.py}, at the
commit in~\cite{baseline}. Exact full-validator control for the independent
scope-fusion supplement.
\end{thebibliography}
\end{document}
